\documentclass[12pt,a4paper]{article}

% Geometry & layout
\usepackage[margin=1in]{geometry}
\setlength{\parskip}{0.5em}

% Fonts & encoding
\usepackage[T1]{fontenc}
\usepackage[expansion=false]{microtype}

% Hyperlinks
\PassOptionsToPackage{hyphens}{url}
\usepackage{hyperref}
\hypersetup{
    colorlinks=true,
    linkcolor=blue,
    urlcolor=blue,
    citecolor=blue,
    breaklinks=true
}

% Tables
\usepackage{booktabs}
\usepackage{longtable}
\usepackage{tabularx}
\usepackage{multirow}
\usepackage{array}
\usepackage{makecell}

% Graphics & plots
\usepackage{graphicx}
\usepackage{tikz}
\usepackage{pgfplots}
\pgfplotsset{compat=1.18}
\usepackage{pgfplotstable}

% Math
\usepackage{amsmath}

% Lists
\usepackage{enumitem}

% Captions & subfigures
\usepackage{caption}
\usepackage{subcaption}

% Code listings
\usepackage{listings}
\usepackage{xcolor}

\lstset{
    basicstyle=\ttfamily\footnotesize,
    breaklines=true,
    frame=single,
    numbers=none,
    showstringspaces=false,
    keywordstyle=\color{blue},
    commentstyle=\color{gray},
    backgroundcolor=\color{gray!10}
}

% Headers/formatting
\usepackage{fancyhdr}
\usepackage{titlesec}

% Floats
\usepackage{float}

% Custom spacing tweak
\emergencystretch=3em

% Path command
\newcommand{\pth}[1]{\path{#1}}

\hypersetup{
    colorlinks=true,
    linkcolor=blue,
    urlcolor=blue,
    citecolor=blue,
    breaklinks=true
}

\lstset{
    basicstyle=\ttfamily\footnotesize,
    breaklines=true,
    frame=single,
    numbers=none,
    showstringspaces=false,
    keywordstyle=\color{blue},
    commentstyle=\color{gray},
    backgroundcolor=\color{gray!10}
}

\begin{document}

\begin{titlepage}
    \centering

    \vspace*{1.5cm}

    {\Large \textbf{Indian Institute of Technology Kanpur}}\\[0.6em]
    {\normalsize Department of Computer Science and Engineering}

    \vspace{2.5cm}

    {\large \textbf{CS614: Linux Kernel Programming}}\\[0.6em]
    {\normalsize \textit{Final Report}}

    \vspace{2.5cm}

    {\LARGE \textbf{Analysis, Benchmarking, and Optimization of}}\\[0.4em]
    {\LARGE \textbf{ext4/JBD2 Consistency Mechanism on Linux 6.1.4}}

    \vspace{3cm}

    {\large \textbf{Team Members}}\\[0.8em]
    \begin{tabular}{ll}
        \toprule
        \textbf{Name} & \textbf{Roll Number} \\
        \midrule
        Milan Roy        & 241110042 \\
        Sahil Basia      & 241110061 \\
        Shrey Sharma     & 251110068 \\
        Suyamoon Pathak  & 241110091 \\
        \bottomrule
    \end{tabular}

    \vfill

    {\normalsize April 2026}

    \vspace{1cm}

\end{titlepage}

\begin{abstract}
This report describes the work of the \textbf{MTech Rocks} group on
the ext4/JBD2 consistency mechanism on Linux~6.1.4. We benchmarked
the baseline cost of each ext4 journaling mode and identified where
the journal hurts performance most; this characterization provided
clear directions for further optimization. We then addressed two
fast-commit fallback paths with small kernel patches:
\texttt{EXT4\_FC\_REASON\_XATTR} (108 lines, 64$\times$ fewer JBD2
full commits on a setxattr workload, 27\% wall-time gain on VM and
48\% on bare-metal) and
\texttt{EXT4\_FC\_REASON\_FALLOC\_RANGE} (46 lines, 62.5$\times$
fewer commits on collapse/insert range, 23--26\% improvement on VM
and 42--44\% on bare-metal).

We further developed a kernel module using \texttt{kretprobes} to
measure that $\sim$78\% of fsync time is spent inside
\texttt{jbd2\_log\_wait\_commit}, and replaced this function with a
lockless CAS-based version, achieving improvements of +40\%, +54\%,
and +7\% IOPS at 16 threads across three workloads, without
regression in streaming performance.

Finally, we designed a kernel module that replaces JBD2's strict
ring-buffer journal allocator with a free-block bitmap, reducing
P99 latency on a WAL workload from 0.41~ms to 0.05~ms in
\texttt{ordered} mode while maintaining the same write
amplification factor.
\end{abstract}

\newpage
\tableofcontents
\newpage

% =====================================================================
\section{Introduction}
% =====================================================================

Modern filesystems must guarantee consistency in the face of unexpected system crashes
or power failures. Without such guarantees, an interrupted write sequence can leave the
filesystem in a partially updated state, corrupting metadata structures such as inodes,
directory entries, and block allocation bitmaps. Recovery from such corruption is
expensive, often requiring a full filesystem check (\texttt{fsck}), and may result in
permanent data loss.

To address this, ext4 --- the default filesystem on most Linux distributions --- employs
a journaling layer called \textbf{JBD2} (Journaling Block Device 2). Rather than writing
updates directly to their final on-disk locations, JBD2 first records changes atomically
to a dedicated circular log called the \textit{journal}. On crash recovery, the kernel
replays any committed but not yet checkpointed transactions from the journal, restoring
the filesystem to a consistent state without a full scan. 

JBD2 organises writes into \textit{transactions}: a group of modifications that are
committed atomically. Each transaction passes through a well-defined pipeline ---
accumulating updates, locking for commit, writing descriptor and data blocks to the
journal, issuing a commit block, and finally checkpointing dirty blocks to their home
locations. Crucially, a transaction is not considered durable until a commit block is
flushed to stable storage, enforcing the write-ahead property.

Ext4 exposes three journaling modes, each offering a different trade-off between
consistency and performance:

\begin{itemize}
    \item \textbf{Journal mode}: Both file data and metadata are written to the journal
    before being committed to their final locations. This provides the strongest
    consistency guarantees --- even file data is recoverable after a crash --- but incurs
    the highest overhead due to the double-write of all data blocks.

    \item \textbf{Ordered mode} (default): Only metadata is journaled. However, ext4
    enforces a strict ordering constraint: data blocks must be flushed to their final
    locations \emph{before} the corresponding metadata transaction is committed to the
    journal. This prevents stale data from being exposed after a crash, while avoiding
    the cost of journalling data blocks directly.

    \item \textbf{Writeback mode}: Only metadata is journaled, with no ordering
    constraint between data and metadata writes. This mode offers the best write
    throughput but the weakest consistency guarantee: after a crash, a committed metadata
    update may point to a data block that was never fully written, potentially exposing
    stale or garbage data.
\end{itemize}

% =====================================================================
\section{Workload Characterization}
% =====================================================================
\section{Workload Characterisation}
% =====================================================================

\subsection{Sync-Heavy Small Writes Workload}
\subsubsection{Workload Description}
This workload emulates the write-ahead logging (WAL) pattern typical of database
systems, in which each log record must be made durable before the next write can
proceed. Using \texttt{fio}, a single thread writes a total of 100\,MB to a file in
fixed-size blocks of 4, 8, 16, 32, or 64\,KB. An \texttt{fsync} is issued after every
block write, forcing a JBD2 transaction commit to complete before the subsequent write
begins. The benchmark is executed across all three ext4 journaling modes ---
\texttt{ordered}, \texttt{journal}, and \texttt{writeback} --- together with a
\texttt{nojournal} baseline, enabling a controlled comparison of journaling overhead
and write amplification across block sizes. Each configuration is repeated five times
and the results are averaged to mitigate measurement variability.

\begin{table}[H]
\centering
\caption{Throughput, Latency, and I/O Breakdown across Journaling Modes and Block Sizes}
\label{tab:throughput}
\renewcommand{\arraystretch}{1.1}
\small
\resizebox{\textwidth}{!}{%
\begin{tabular}{llrrrrrrr}
\toprule
\textbf{Mode} & \textbf{BS}
  & \textbf{IOPS}
  & \textbf{BW (MB/s)}
  & \textbf{Runtime (s)}
  & \textbf{Avg lat.\ (ms)}
  & \textbf{write() (µs)}
  & \textbf{fsync() (µs)}
  & \textbf{fsync \%} \\
\midrule
\multirow{5}{*}{ordered}
  & 4K  & $185 \pm 4$   & $0.72 \pm 0.01$ & $138.05 \pm 2.70$  & $5.39 \pm 0.11$ & $20 \pm 0$ & $5369 \pm 105$ & 99.6\% \\
  & 8K  & $180 \pm 5$   & $1.40 \pm 0.04$ & $71.24  \pm 2.20$  & $5.56 \pm 0.17$ & $22 \pm 0$ & $5541 \pm 172$ & 99.6\% \\
  & 16K & $174 \pm 4$   & $2.71 \pm 0.06$ & $36.85  \pm 0.83$  & $5.76 \pm 0.13$ & $25 \pm 0$ & $5731 \pm 130$ & 99.6\% \\
  & 32K & $173 \pm 2$   & $5.39 \pm 0.06$ & $18.55  \pm 0.19$  & $5.80 \pm 0.06$ & $28 \pm 0$ & $5767 \pm 61$  & 99.5\% \\
  & 64K & $170 \pm 3$   & $10.62\pm 0.21$ & $9.41   \pm 0.19$  & $5.89 \pm 0.12$ & $39 \pm 0$ & $5845 \pm 119$ & 99.3\% \\
\midrule
\multirow{5}{*}{journal}
  & 4K  & $179 \pm 3$   & $0.70 \pm 0.01$ & $142.78 \pm 2.63$  & $5.57 \pm 0.10$ & $23 \pm 0$ & $5551 \pm 103$ & 99.6\% \\
  & 8K  & $169 \pm 2$   & $1.32 \pm 0.01$ & $75.63  \pm 0.71$  & $5.91 \pm 0.06$ & $26 \pm 0$ & $5880 \pm 55$  & 99.6\% \\
  & 16K & $166 \pm 3$   & $2.59 \pm 0.04$ & $38.65  \pm 0.60$  & $6.04 \pm 0.09$ & $31 \pm 0$ & $6006 \pm 93$  & 99.5\% \\
  & 32K & $165 \pm 3$   & $5.17 \pm 0.11$ & $19.36  \pm 0.41$  & $6.05 \pm 0.13$ & $40 \pm 0$ & $6007 \pm 129$ & 99.3\% \\
  & 64K & $154 \pm 7$   & $9.62 \pm 0.47$ & $10.41  \pm 0.51$  & $6.51 \pm 0.32$ & $59 \pm 1$ & $6447 \pm 320$ & 99.1\% \\
\midrule
\multirow{5}{*}{writeback}
  & 4K  & $185 \pm 2$   & $0.72 \pm 0.01$ & $138.21 \pm 1.33$  & $5.40 \pm 0.05$ & $20 \pm 0$ & $5376 \pm 52$  & 99.6\% \\
  & 8K  & $181 \pm 1$   & $1.41 \pm 0.01$ & $70.88  \pm 0.27$  & $5.53 \pm 0.02$ & $22 \pm 0$ & $5513 \pm 21$  & 99.6\% \\
  & 16K & $170 \pm 1$   & $2.66 \pm 0.02$ & $37.58  \pm 0.28$  & $5.87 \pm 0.04$ & $25 \pm 0$ & $5844 \pm 44$  & 99.6\% \\
  & 32K & $160 \pm 7$   & $5.00 \pm 0.22$ & $20.05  \pm 0.88$  & $6.26 \pm 0.27$ & $29 \pm 0$ & $6234 \pm 275$ & 99.5\% \\
  & 64K & $159 \pm 6$   & $9.93 \pm 0.40$ & $10.08  \pm 0.43$  & $6.30 \pm 0.27$ & $40 \pm 1$ & $6262 \pm 268$ & 99.4\% \\
\midrule
\multirow{5}{*}{nojournal}
  & 4K  & $381 \pm 3$   & $1.49 \pm 0.01$ & $67.17  \pm 0.56$  & $2.62 \pm 0.02$ & $15 \pm 0$ & $2604 \pm 22$  & 99.4\% \\
  & 8K  & $362 \pm 5$   & $2.83 \pm 0.04$ & $35.38  \pm 0.50$  & $2.76 \pm 0.04$ & $17 \pm 0$ & $2743 \pm 39$  & 99.4\% \\
  & 16K & $296 \pm 5$   & $4.62 \pm 0.08$ & $21.63  \pm 0.37$  & $3.38 \pm 0.06$ & $20 \pm 0$ & $3356 \pm 59$  & 99.4\% \\
  & 32K & $307 \pm 8$   & $9.60 \pm 0.26$ & $10.43  \pm 0.28$  & $3.26 \pm 0.09$ & $24 \pm 0$ & $3230 \pm 86$  & 99.3\% \\
  & 64K & $269 \pm 13$  & $16.82\pm 0.79$ & $5.95   \pm 0.29$  & $3.72 \pm 0.18$ & $34 \pm 0$ & $3684 \pm 181$ & 99.1\% \\
\bottomrule
\end{tabular}%
}
\end{table}

\subsubsection{Throughput and Latency}
Prior to benchmarking, we expected \texttt{ordered} mode to exhibit slightly lower IOPS
than \texttt{writeback} mode, owing to the additional ordering constraint that data
blocks must be flushed to disk before their corresponding metadata is committed to the
journal. We further expected \texttt{journal} mode to incur substantially lower IOPS
than either of the other two modes, given the overhead of writing file data into the
journal in addition to metadata. Table~\ref{tab:throughput} shows neither expectation
to be borne out: IOPS across all three journaling modes are remarkably similar, whilst
\texttt{nojournal} achieves roughly twice the IOPS of any journaled mode.

The near-equivalence among the three journaled modes follows from the fact that all
three must traverse the same JBD2 commit pipeline on every \texttt{fsync}: an NVMe
cache flush to persist both data and journal log blocks to their on-disk locations,
followed by a direct write of the commit block to stable storage. The additional data
written into the journal by \texttt{journal} mode contributes only a marginal cost at
smaller block sizes, insufficient to produce a measurable IOPS gap. \texttt{nojournal}
mode, by contrast, incurs only the first of these two costs --- a single NVMe cache
flush --- and writes no journal entries whatsoever, halving the per-\texttt{fsync}
barrier count and eliminating journal write amplification entirely. This accounts for
why \texttt{nojournal} achieves approximately twice the IOPS of the journaled modes,
particularly at smaller block sizes. As block size increases, IOPS decline modestly
across all modes, reflecting the growing per-operation data transfer cost beginning to
contribute alongside the fixed flush overhead.

Although per-I/O latency and IOPS remain roughly constant across block sizes for a
given mode, total benchmark time roughly halves with each doubling of block size.
Doubling the block size doubles the data written per \texttt{fsync} call, so the
total number of \texttt{fsync} calls required is halved proportionally. Since the
per-\texttt{fsync} cost is dominated by the fixed NVMe flush latency, total time falls
in proportion. Notably, total time is consistent across all journaling modes at a given
block size, including \texttt{journal} mode, for the same reason: the additional data
journaled by \texttt{journal} mode does not meaningfully increase the per-flush cost.

\subsubsection{Impact of fsync}
As shown in Table~\ref{tab:throughput}, the \texttt{write()} system call contributes at
most 0.9\% of total per-I/O latency across all modes and block sizes, since it
merely copies data into the page cache and returns immediately. The \texttt{write()}
latency increases slightly with block size --- most visibly in \texttt{journal} mode
--- but does not double when the block size doubles, indicating a fixed component.

The \texttt{fsync()} latency, by contrast, is two orders of magnitude higher and
subsumes the cost of the JBD2 commit pipeline and the NVMe cache flush.
\texttt{fsync()} similarly does not double with block size, reflecting the large fixed
component of NVMe flush latency that is independent of the volume of data being
flushed, at least at smaller block sizes. Consequently, \texttt{fsync()} dominates
total per-I/O latency in every configuration, conclusively establishing that the
bottleneck is the mandatory durability guarantee enforced by \texttt{fsync()}, not
the movement of data into the page cache.

\subsubsection{Physical Bytes and Write Amplification}

\begin{table}[H]
\centering
\caption{Write Amplification, Per-Transaction Workload, and JBD2 Commit Pipeline across Journaling Modes and Block Sizes}
\label{tab:waf_txn}
\renewcommand{\arraystretch}{1.1}
\small
\resizebox{\textwidth}{!}{%
\begin{tabular}{llrrrrrrrrr}
\toprule
\textbf{Mode} & \textbf{BS}
  & \thead{\textbf{Physical}\\\textbf{(MB)}}
  & \thead{\textbf{Extra}\\\textbf{Bytes (MB)}}
  & \textbf{WAF}
  & \thead{\textbf{Total}\\\textbf{txns}}
  & \thead{\textbf{Blocks}\\\textbf{logged/txn}}
  & \thead{\textbf{Avg}\\\textbf{Logging}\\\textbf{(ms)}}
  & \thead{\textbf{Avg E2E}\\\textbf{commit}\\\textbf{(ms)}}
  & \thead{\textbf{JBD2}\\\textbf{commit \%}} \\
\midrule
\multirow{5}{*}{ordered}
  & 4K  & $400.92 \pm 1.79$ & $300.92 \pm 1.79$ & $2.0065 \pm 0.0090$ & $25600 \pm 1$ & $3.0 \pm 0.0$ & $5.268 \pm 0.102$ & $5.272 \pm 0.102$ & $97.8 \pm 1.9$\% \\
  & 8K  & $250.08 \pm 0.01$ & $150.08 \pm 0.01$ & $1.6670 \pm 0.0001$ & $12801 \pm 1$ & $3.0 \pm 0.0$ & $5.437 \pm 0.172$ & $5.441 \pm 0.172$ & $97.8 \pm 3.1$\% \\
  & 16K & $175.06 \pm 0.01$ & $ 75.06 \pm 0.01$ & $1.4002 \pm 0.0000$ & $ 6401 \pm 0$ & $3.0 \pm 0.0$ & $5.621 \pm 0.129$ & $5.625 \pm 0.129$ & $97.7 \pm 2.2$\% \\
  & 32K & $137.56 \pm 0.00$ & $ 37.56 \pm 0.00$ & $1.2225 \pm 0.0000$ & $ 3201 \pm 0$ & $3.0 \pm 0.0$ & $5.651 \pm 0.065$ & $5.655 \pm 0.065$ & $97.6 \pm 1.1$\% \\
  & 64K & $118.81 \pm 0.00$ & $  18.81 \pm 0.00$ & $1.1180 \pm 0.0000$ & $ 1601 \pm 0$ & $3.0 \pm 0.0$ & $5.718 \pm 0.119$ & $5.722 \pm 0.119$ & $97.3 \pm 2.0$\% \\
\midrule
\multirow{5}{*}{journal}
  & 4K  & $500.14 \pm 0.00$ & $400.14 \pm 0.00$ & $2.5031 \pm 0.0008$ & $25600 \pm 0$ & $ 4.0 \pm 0.0$ & $5.460 \pm 0.100$ & $5.467 \pm 0.100$ & $98.0 \pm 1.8$\% \\
  & 8K  & $350.10 \pm 0.00$ & $250.10 \pm 0.00$ & $2.3336 \pm 0.0001$ & $12800 \pm 0$ & $ 5.0 \pm 0.0$ & $5.778 \pm 0.051$ & $5.785 \pm 0.051$ & $97.9 \pm 0.9$\% \\
  & 16K & $275.07 \pm 0.00$ & $175.07 \pm 0.00$ & $2.2001 \pm 0.0000$ & $ 6400 \pm 0$ & $ 7.0 \pm 0.0$ & $5.896 \pm 0.093$ & $5.903 \pm 0.093$ & $97.7 \pm 1.5$\% \\
  & 32K & $237.54 \pm 0.00$ & $137.54 \pm 0.00$ & $2.1111 \pm 0.0000$ & $ 3200 \pm 0$ & $11.0 \pm 0.0$ & $5.889 \pm 0.130$ & $5.897 \pm 0.130$ & $97.5 \pm 2.1$\% \\
  & 64K & $218.76 \pm 0.00$ & $118.76 \pm 0.00$ & $2.0584 \pm 0.0000$ & $ 1600 \pm 0$ & $19.0 \pm 0.0$ & $6.315 \pm 0.318$ & $6.323 \pm 0.318$ & $97.2 \pm 4.9$\% \\
\midrule
\multirow{5}{*}{writeback}
  & 4K  & $400.11 \pm 0.01$ & $300.11 \pm 0.01$ & $2.0024 \pm 0.0006$ & $25600 \pm 0$ & $3.0 \pm 0.0$ & $5.275 \pm 0.051$ & $5.279 \pm 0.051$ & $97.8 \pm 1.0$\% \\
  & 8K  & $250.09 \pm 0.00$ & $150.09 \pm 0.00$ & $1.6670 \pm 0.0001$ & $12801 \pm 0$ & $3.0 \pm 0.0$ & $5.406 \pm 0.023$ & $5.410 \pm 0.023$ & $97.7 \pm 0.4$\% \\
  & 16K & $175.06 \pm 0.01$ & $ 75.06 \pm 0.01$ & $1.4002 \pm 0.0001$ & $ 6401 \pm 1$ & $3.0 \pm 0.0$ & $5.731 \pm 0.044$ & $5.735 \pm 0.044$ & $97.7 \pm 0.8$\% \\
  & 32K & $137.56 \pm 0.00$ & $ 37.56 \pm 0.00$ & $1.2225 \pm 0.0000$ & $ 3201 \pm 0$ & $3.0 \pm 0.0$ & $6.112 \pm 0.272$ & $6.116 \pm 0.272$ & $97.7 \pm 4.3$\% \\
  & 64K & $118.81 \pm 0.00$ & $  18.81 \pm 0.00$ & $1.1180 \pm 0.0000$ & $ 1601 \pm 0$ & $3.0 \pm 0.0$ & $6.130 \pm 0.265$ & $6.135 \pm 0.265$ & $97.4 \pm 4.2$\% \\
\midrule
\multirow{5}{*}{nojournal}
  & 4K  & $199.81 \pm 0.06$ & $ 99.81 \pm 0.06$ & $1.0000 \pm 0.0004$ & --- & --- & --- & --- & --- \\
  & 8K  & $150.02 \pm 0.01$ & $ 50.02 \pm 0.01$ & $1.0000 \pm 0.0001$ & --- & --- & --- & --- & --- \\
  & 16K & $125.02 \pm 0.00$ & $ 25.02 \pm 0.00$ & $1.0000 \pm 0.0000$ & --- & --- & --- & --- & --- \\
  & 32K & $112.52 \pm 0.00$ & $ 12.52 \pm 0.00$ & $1.0000 \pm 0.0000$ & --- & --- & --- & --- & --- \\
  & 64K & $106.27 \pm 0.00$ & $  6.27 \pm 0.00$ & $1.0000 \pm 0.0000$ & --- & --- & --- & --- & --- \\
\bottomrule
\end{tabular}%
}
\end{table}

The Physical and Extra Bytes columns in Table~\ref{tab:waf_txn} capture the disk-write
overhead beyond the 100\,MB logical payload, as measured by \texttt{blktrace}. An
important detail when interpreting these values is that disk writes are granular to
4\,KB: even a few bytes of metadata require a full 4\,KB block, inflating the observed
overhead accordingly.

Across all modes, extra bytes halve as block size doubles, since the transaction count
halves proportionally and each transaction carries a fixed metadata footprint. In
\texttt{journal} mode this halving is partially masked by an additional
${\approx}100$\,MB term arising from the data blocks being written twice into the
journal, but removing this constant reveals the same underlying pattern.

For \texttt{ordered} and \texttt{writeback} modes, which log exactly three blocks per
transaction --- one descriptor, one metadata, and one commit block --- the extra bytes
can be verified analytically:
\begin{equation}
    3 \times 25{,}600 \text{ transactions} \times 4\,\text{KB} \approx 300\,\text{MB}
\end{equation}
This matches the observed value at 4\,KB, and the same halving applies at larger block
sizes. An analogous calculation for \texttt{journal} mode must additionally account for
the data blocks logged per transaction, which scale linearly with block size.

WAF (Table~\ref{tab:waf_txn}) quantifies how many times more data is written to disk
relative to the \texttt{nojournal} baseline. For \texttt{ordered} and
\texttt{writeback}, WAF decreases from $2.00\times$ at 4\,KB to $1.12\times$ at
64\,KB, as the fixed three-block journal overhead is progressively diluted by the
growing data payload. \texttt{journal} mode converges far more slowly, from
$2.50\times$ to $2.06\times$, because its data-journaling cost scales with block size
and cannot be diluted in the same manner.

\subsubsection{JBD2 Commit Percentage of Benchmark Time}
As shown in Table~\ref{tab:waf_txn}, the average logging time --- the time spent
writing the journal transaction to the NVMe device --- accounts for virtually all of
the end-to-end commit time across every configuration. JBD2 commit activity
consequently dominates total benchmark wall-clock time, a direct consequence of the
\texttt{fsync}-after-every-write pattern that chains each write to the completion of
the preceding commit. With the application blocked almost entirely on synchronous I/O,
there is no opportunity for any form of parallelism or batching.

\subsubsection{Conclusion}
The sync-heavy WAL workload establishes that, under a per-write \texttt{fsync}
discipline, the dominant cost is the NVMe cache flush that \texttt{fsync()} must issue
on every commit, not the movement of data or metadata within the journaling layer. All
three journaled modes exhibit near-identical IOPS because they share the same commit
pipeline; the additional data written by \texttt{journal} mode is insufficient to
widen this gap at the block sizes examined. \texttt{nojournal} mode achieves
approximately twice the IOPS by eliminating the journal write entirely and incurring
only a single flush per \texttt{fsync}. WAF decreases with block size for
\texttt{ordered} and \texttt{writeback} as the fixed three-block overhead is diluted by
larger payloads, whereas \texttt{journal} mode's WAF converges only slowly because its
data-journaling cost grows with block size. 

% =====================================================================
\subsection{Sequential Write Workload}

\subsubsection{Workload Description}
This workload measures the raw sequential write performance of each ext4 journaling
mode without the durability constraint of a per-write \texttt{fsync}, emulating the
log-writing behaviour common in database and application logging scenarios.
\texttt{fio} is configured as a single-threaded workload writing a total of 10\,GB of
data to a file in fixed-size blocks of 128\,KB, 256\,KB, 512\,KB, and 1\,MB. Unlike
the sync-heavy WAL workload, no \texttt{fsync} is issued between writes, allowing the
kernel to freely batch dirty pages and pipeline writes to disk. JBD2 commits are
therefore not driven by the application but are instead triggered by the kernel's
background commit timer, which fires every 5\,seconds to lock the current running
transaction and flush it to disk. The benchmark is executed across all three ext4
journaling modes --- \texttt{ordered}, \texttt{journal}, \texttt{writeback}, and
\texttt{nojournal} --- to isolate the cost of each mode's journaling policy under
conditions where the kernel is free to amortise overhead across large batches of
writes. Each configuration is repeated five times and the results are averaged to
mitigate measurement variability.

\subsubsection{Throughput and Latency}
\begin{table}[H]
\centering
\caption{Throughput, Latency, and I/O Breakdown across Journaling Modes and Block Sizes (Sequential Write)}
\label{tab:seq_throughput}
\renewcommand{\arraystretch}{1.1}
\small
\resizebox{\textwidth}{!}{%
\begin{tabular}{llrrrrrr}
\toprule
\textbf{Mode} & \textbf{BS}
  & \textbf{IOPS}
  & \textbf{BW (MB/s)}
  & \textbf{Runtime (s)}
  & \textbf{Avg (ms)}
  & \textbf{write() (µs)} \\
\midrule
\multirow{4}{*}{ordered}
  & 128K & $1580 \pm 772$  & $197.54 \pm 96.47$ & $65.35  \pm 42.10$ & $0.80 \pm 0.51$ & $796  \pm 514$ \\
  & 256K & $693  \pm 160$  & $173.17 \pm 40.01$ & $61.14  \pm 13.11$ & $1.49 \pm 0.32$ & $1488 \pm 319$ \\
  & 512K & $491  \pm 81$   & $245.67 \pm 40.48$ & $42.49  \pm  7.41$ & $2.07 \pm 0.36$ & $2066 \pm 361$ \\
  & 1M   & $101  \pm 22$   & $100.94 \pm 22.16$ & $105.32 \pm 26.46$ & $10.28\pm 2.58$ & $10251\pm 2584$ \\
\midrule
\multirow{4}{*}{journal}
  & 128K & $290  \pm 7$    & $36.22  \pm  0.93$ & $282.83 \pm  7.26$ & $3.45 \pm 0.09$ & $3451 \pm 89$  \\
  & 256K & $147  \pm 4$    & $36.86  \pm  0.92$ & $277.89 \pm  7.04$ & $6.78 \pm 0.17$ & $6781 \pm 172$ \\
  & 512K & $74   \pm 4$    & $37.17  \pm  2.12$ & $276.08 \pm 16.28$ & $13.48\pm 0.80$ & $13473\pm 796$ \\
  & 1M   & $28   \pm 0$    & $27.96  \pm  0.20$ & $366.30 \pm  2.69$ & $35.77\pm 0.26$ & $35750\pm 262$ \\
\midrule
\multirow{4}{*}{writeback}
  & 128K & $993  \pm 44$   & $124.15 \pm  5.55$ & $82.59  \pm  3.61$ & $1.01 \pm 0.04$ & $1006 \pm 44$  \\
  & 256K & $520  \pm 8$    & $130.03 \pm  1.92$ & $78.76  \pm  1.15$ & $1.92 \pm 0.03$ & $1919 \pm 28$  \\
  & 512K & $270  \pm 21$   & $134.80 \pm 10.46$ & $76.26  \pm  5.78$ & $3.72 \pm 0.28$ & $3712 \pm 284$ \\
  & 1M   & $82   \pm 1$    & $81.73  \pm  0.93$ & $125.30 \pm  1.42$ & $12.23\pm 0.14$ & $12201\pm 139$ \\
\midrule
\multirow{4}{*}{nojournal}
  & 128K & $1362 \pm 529$  & $170.24 \pm 66.07$ & $67.92  \pm 30.58$ & $0.83 \pm 0.37$ & $826  \pm 373$ \\
  & 256K & $714  \pm 67$   & $178.47 \pm 16.69$ & $57.72  \pm  5.60$ & $1.41 \pm 0.14$ & $1405 \pm 137$ \\
  & 512K & $413  \pm 106$  & $206.60 \pm 52.97$ & $51.59  \pm 11.93$ & $2.52 \pm 0.58$ & $2508 \pm 582$ \\
  & 1M   & $96   \pm 2$    & $95.77  \pm  2.10$ & $106.95 \pm  2.33$ & $10.44\pm 0.23$ & $10407\pm 232$ \\
\bottomrule
\end{tabular}%
}
\end{table}

The most striking feature of Table~\ref{tab:seq_throughput} is the sharp and persistent
throughput gap between \texttt{journal} mode and the remaining three modes across all
block sizes. Without a per-write \texttt{fsync} serialising every commit, the kernel
is free to batch dirty pages and commit journal transactions in the background without
blocking the application thread: page cache writeback and JBD2 transaction commits
proceed concurrently with the benchmarked workload, and since no explicit flush command
is issued to the NVMe device, the cache flush cost is absorbed in parallel as well.
This removes the dominant bottleneck of the sync-heavy workload and exposes the true
cost of each mode's journaling policy in isolation.

A notable contrast with the sync-heavy workload is that total benchmark time in
Table~\ref{tab:seq_throughput} is largely independent of block size for
\texttt{ordered}, \texttt{writeback}, and \texttt{nojournal}. In the sync-heavy
workload, total time halved with each doubling of block size because each
\texttt{fsync} incurred a fixed NVMe flush cost and larger blocks reduced the number of
\texttt{fsync} calls required. Here, with no per-write \texttt{fsync}, the flush cost
is no longer on the critical path: the workload, page cache writeback, and journal
commits all proceed concurrently, keeping the NVMe device continuously busy regardless
of block size. This parallelism also explains why the IOPS, bandwidth, and
\texttt{write()} latency of \texttt{ordered} and \texttt{writeback} remain close to the
\texttt{nojournal} baseline --- the journaling overhead of these modes is absorbed
entirely in the background and does not delay the application thread.

\subsubsection{Write Amplification and Commit Behaviour}

\begin{table}[H]
\centering
\caption{Write Amplification and JBD2 Commit Behaviour across Journaling Modes and Block Sizes (Sequential Write)}
\label{tab:seq_waf}
\renewcommand{\arraystretch}{1.1}
\small
\resizebox{\textwidth}{!}{%
\begin{tabular}{llrrrrrr}
\toprule
\textbf{Mode} & \textbf{BS}
  & \thead{\textbf{Physical}\\\textbf{(MB)}}
  & \textbf{WAF}
  & \thead{\textbf{Total}\\\textbf{Commits}}
  & \thead{\textbf{Total}\\\textbf{Time (s)}}
  & \thead{\textbf{Commit}\\\textbf{interval (s)}} \\
\midrule
\multirow{4}{*}{ordered}
  & 128K & $10240.94 \pm 0.13$ & $1.0000 \pm 0.0000$ & $16 \pm 8$  & $65.35  \pm 42.10$ & $4.084 \pm 3.331$ \\
  & 256K & $10240.92 \pm 0.03$ & $1.0000 \pm 0.0000$ & $14 \pm 2$  & $61.14  \pm 13.11$ & $4.367 \pm 1.125$ \\
  & 512K & $10240.87 \pm 0.02$ & $1.0000 \pm 0.0000$ & $11 \pm 1$  & $42.49  \pm  7.41$ & $3.863 \pm 0.760$ \\
  & 1M   & $10241.04 \pm 0.08$ & $1.0000 \pm 0.0000$ & $22 \pm 5$  & $105.32 \pm 26.46$ & $4.787 \pm 1.622$ \\
\midrule
\multirow{4}{*}{journal}
  & 128K & $20522.52 \pm 1.22$ & $2.0040 \pm 0.0001$ & $622 \pm 42$ & $282.83 \pm  7.26$ & $0.455 \pm 0.033$ \\
  & 256K & $20523.58 \pm 0.94$ & $2.0041 \pm 0.0001$ & $646 \pm 0$  & $277.89 \pm  7.04$ & $0.430 \pm 0.011$ \\
  & 512K & $20522.95 \pm 0.63$ & $2.0041 \pm 0.0001$ & $646 \pm 0$  & $276.08 \pm 16.28$ & $0.427 \pm 0.025$ \\
  & 1M   & $20523.61 \pm 0.26$ & $2.0041 \pm 0.0001$ & $646 \pm 0$  & $366.30 \pm  2.69$ & $0.567 \pm 0.004$ \\
\midrule
\multirow{4}{*}{writeback}
  & 128K & $10240.99 \pm 0.02$ & $1.0000 \pm 0.0000$ & $18 \pm 2$  & $82.59  \pm  3.61$ & $4.588 \pm 0.548$ \\
  & 256K & $10240.98 \pm 0.01$ & $1.0000 \pm 0.0000$ & $18 \pm 1$  & $78.76  \pm  1.15$ & $4.376 \pm 0.251$ \\
  & 512K & $10240.97 \pm 0.01$ & $1.0000 \pm 0.0000$ & $18 \pm 1$  & $76.26  \pm  5.78$ & $4.237 \pm 0.398$ \\
  & 1M   & $10241.10 \pm 0.02$ & $1.0000 \pm 0.0000$ & $26 \pm 1$  & $125.30 \pm  1.42$ & $4.819 \pm 0.193$ \\
\midrule
\multirow{4}{*}{nojournal}
  & 128K & $10240.56 \pm 0.08$ & $1.0000 \pm 0.0000$ & --- & $67.92  \pm 30.58$ & --- \\
  & 256K & $10240.58 \pm 0.01$ & $1.0000 \pm 0.0000$ & --- & $57.72  \pm  5.60$ & --- \\
  & 512K & $10240.59 \pm 0.05$ & $1.0000 \pm 0.0000$ & --- & $51.59  \pm 11.93$ & --- \\
  & 1M   & $10240.70 \pm 0.24$ & $1.0000 \pm 0.0000$ & --- & $106.95 \pm  2.33$ & --- \\
\bottomrule
\end{tabular}%
}
\end{table}

Table~\ref{tab:seq_waf} consolidates the physical write volume, WAF, and commit
behaviour for each mode. Unlike the sync-heavy workload, where physical bytes written
decreased as block size doubled due to reduced per-transaction overhead, here the
physical byte count is essentially constant across all block sizes within each mode.
\texttt{ordered}, \texttt{writeback}, and \texttt{nojournal} each produce approximately
10.240\,GB of physical writes, whilst \texttt{journal} mode consistently produces
around 20.522\,GB. This block-size independence arises because JBD2 commits are driven
by the background timer rather than by per-write \texttt{fsync} calls, leaving the
physical write volume determined entirely by the data volume and the mode's journaling
policy rather than by write granularity.

The WAF results reflect this directly. \texttt{ordered} and \texttt{writeback} achieve
a WAF of effectively $1.0000\times$ across all block sizes, identical to
\texttt{nojournal}. With the background timer batching thousands of writes per
transaction --- only 14--28 total commits are issued across the entire 10\,GB run ---
the fixed per-transaction overhead of descriptor, metadata, and commit blocks is
amortised to negligible levels. This contrasts sharply with the sync-heavy workload,
where \texttt{ordered} mode carried a WAF of $2.00\times$ at 4\,KB because every
write forced its own individual commit. \texttt{journal} mode, by contrast, maintains a
WAF of exactly $2.004\times$ across every block size with no variation. Because every
data block must be written into the journal in addition to its final disk location, the
extra write cost is strictly proportional to the volume of data written rather than to
the number of transactions, and no amount of JBD2 batching can reduce this inherent
$2\times$ factor.

The commit interval column makes the underlying commit mechanism explicit.
\texttt{ordered} and \texttt{writeback} show commit intervals of 4.0--4.8\,s across
all block sizes, consistent with the kernel's 5\,second background commit timer firing
with no other trigger. \texttt{journal} mode, by contrast, exhibits a commit interval
of only 0.427--0.567\,s, roughly $20\times$ more frequent than the timer period. This
is because each transaction in \texttt{journal} mode must accommodate both data and
metadata blocks, and the journal's fixed maximum transaction capacity is exhausted long
before the 5\,second timer fires. Once the transaction fills its block allocation, JBD2
forces an immediate commit regardless of elapsed time, resulting in 622--646 commits
over the course of the benchmark compared to only 11--26 for the other journaled modes.

\subsubsection{Conclusion}
The sequential write workload exposes a fundamental asymmetry in how the four journaling
modes respond to the removal of the per-write durability constraint. \texttt{ordered}
and \texttt{writeback} exploit JBD2 transaction batching to amortise their
per-transaction journal overhead to negligible levels, achieving near-identical
bandwidth to \texttt{nojournal} and a WAF of effectively $1.0000\times$ across all
block sizes. \texttt{journal} mode is structurally unable to benefit from this
batching: its $2\times$ data amplification is proportional to the volume of data
written rather than to the number of transactions, resulting in a constant WAF of
$2.004\times$, a commit interval $20\times$ shorter than the background timer period,
and a persistent throughput penalty regardless of block size or transaction depth. For
workloads that do not require per-write durability, \texttt{journal} mode should
therefore be avoided, as its throughput penalty is inherent to its data-journaling
policy and cannot be mitigated by tuning block size or commit interval.

% =====================================================================
\subsection{Combined Workload}

The two preceding workloads represented opposite extremes of kernel commit behaviour:
one in which a JBD2 commit is forced after every write, and one in which commit timing
is delegated entirely to the kernel's background timer. This workload places both
patterns under simultaneous pressure, running a sync-heavy and a sequential-write
benchmark concurrently so that their interactions can be observed under realistic
mixed-pressure conditions. The sync-heavy component issues 8\,KB writes each followed
by an \texttt{fsync}, whilst the sequential-write component operates with a 1\,MB
block size. Unlike the earlier workloads, which were bounded by a fixed data volume,
both components here run for a fixed duration of two minutes. The benchmark is executed
across all four ext4 journaling modes to isolate the overhead attributable to each
mode's journaling policy. Every configuration is repeated five times and the results
are averaged to mitigate measurement variability.

\begin{table}[H]
\centering
\caption{Combined Throughput and Average Latency for WAL and Sequential Workloads}
\label{tab:combined}
\renewcommand{\arraystretch}{1.1}
\small
\resizebox{\textwidth}{!}{%
\begin{tabular}{llrrrr}
\toprule
\textbf{Workload} & \textbf{Mode}
  & \textbf{IOPS}
  & \textbf{BW (MB/s)}
  & \textbf{Data (GB)}
  & \textbf{Avg lat.\ (ms)} \\
\midrule
\multirow{4}{*}{SYNC}
  & ordered   & $4 \pm 1$   & $0.03 \pm 0.01$ & $0.00 \pm 0.00$ & $252.97 \pm 79.97$  \\
  & journal   & $4 \pm 0$   & $0.03 \pm 0.00$ & $0.00 \pm 0.00$ & $228.39 \pm 18.54$  \\
  & writeback & $2 \pm 1$   & $0.02 \pm 0.01$ & $0.00 \pm 0.00$ & $493.88 \pm 254.09$ \\
  & nojournal & $8 \pm 5$   & $0.06 \pm 0.04$ & $0.01 \pm 0.00$ & $165.53 \pm 98.02$  \\
\midrule
\multirow{4}{*}{SEQ}
  & ordered   & $107 \pm 11$ & $106.74 \pm 11.07$ & $12.52 \pm 1.30$ & $8.83 \pm 0.59$  \\
  & journal   & $34 \pm 2$   & $34.29 \pm 2.12$   & $4.03 \pm 0.25$  & $29.25 \pm 1.69$ \\
  & writeback & $109 \pm 17$ & $108.61 \pm 16.58$ & $12.83 \pm 1.93$ & $9.08 \pm 1.68$  \\
  & nojournal & $96 \pm 13$  & $95.70 \pm 13.44$  & $11.33 \pm 1.43$ & $10.07 \pm 1.60$ \\
\bottomrule
\end{tabular}%
}
\end{table}

\subsubsection{Throughput and Latency}
Comparing Table~\ref{tab:combined} with Table~\ref{tab:throughput}, the WAL component
suffers a severe degradation under concurrent load. IOPS collapse from the
${\sim}170$--$381$ range observed in isolation to just $2$--$8$, a reduction of
95--99\% across all modes. Bandwidth falls correspondingly from up to ${\sim}16.8$\,MB/s
to below $0.06$\,MB/s, and average latency rises dramatically: in \texttt{ordered}
mode, for instance, it increases from ${\sim}5.4$\,ms to ${\sim}253$\,ms, a
$46\times$ inflation. Comparable degradation is observed across all other journaling
modes.

The mechanism is a system-wide coupling introduced by \texttt{fsync}. When the WAL
component calls \texttt{fsync}, the kernel does not flush only the pages belonging to
the WAL file: it locks the current JBD2 transaction and forces a commit of all
outstanding dirty data from every writer sharing the filesystem, including the large
dirty page footprint accumulated by the concurrently running sequential workload. The
\texttt{fsync} call therefore cannot return until this much larger volume of data has
been written back and journaled, inflating each individual \texttt{fsync} latency by
roughly two orders of magnitude relative to its isolated value. Because the WAL thread
is blocked for the entirety of each \texttt{fsync}, its effective throughput collapses.

The sequential workload, by contrast, shows only moderate changes relative to
Table~\ref{tab:seq_throughput}. IOPS and bandwidth are marginally higher in some modes
--- approximately 5--15\% for \texttt{ordered} and \texttt{writeback} --- whilst
average latency decreases slightly. This modest improvement arises because the WAL
component's periodic \texttt{fsync} calls trigger forced JBD2 commits that writeback
the sequential workload's dirty pages as a side effect, reducing the residual writeback
work that would otherwise fall to the background flusher and slightly lowering the
\texttt{write()} latency seen by the sequential thread.

\subsubsection{Conclusion}
The combined workload reveals a cross-workload interference effect that is absent in
either workload individually: a single \texttt{fsync}-issuing thread can severely
degrade its own throughput by inadvertently taking on the writeback cost of all
co-resident writers. The WAL component's \texttt{fsync} latency inflates by up to
$46\times$ relative to its isolated value because each commit is forced to drain the
full dirty page set of the concurrently running sequential workload. This interference
is an intrinsic property of the JBD2 transaction model, in which a commit is a
filesystem-wide operation rather than a per-file one, and its severity scales with the
volume of dirty data accumulated by other writers at the time \texttt{fsync} is called.
The sequential workload is largely unaffected and even benefits slightly, as the
WAL-driven commits serve as additional writeback opportunities that marginally reduce
its own background flusher load. These results highlight that, in multi-workload
environments, the performance of any \texttt{fsync}-dependent application is sensitive
not only to its own I/O pattern but to the aggregate dirty page pressure generated by
all processes sharing the same filesystem.


% =====================================================================
\section{Optimization Candidates: Overview}
% =====================================================================

The measurements in the previous section pointed at the commit path. The natural
question was: which specific commit can we make cheaper or skip
entirely?

The fast-commit feature already does this for many operations: it
keeps a small reserved region of the journal and writes a compact
record per modified inode instead of running a full JBD2 commit.
But fast commit gives up on operations it does not know how to log.
When ext4 hits one of those it calls
\texttt{ext4\_fc\_mark\_ineligible(\_,REASON,\_)} and falls back to
the full commit path. The original FastCommit
paper~\cite{atc24-fc} lists ten such reasons, including
\texttt{XATTR}, \texttt{FALLOC\_RANGE}, \texttt{CROSS\_RENAME},
and \texttt{RENAME\_DIR}. Two of those are what the patches in
Sections~6 and~7 close.

Four candidates were tried in order; only the last two worked.
Section~5 covers C1 and C2 in detail because the lessons they
taught shaped how C3 and C4 were chosen.

\subsection{Candidate selection methodology}

We did not start with fast-commit gaps.  We first code-reviewed the
entire \texttt{fs/jbd2/} and \texttt{fs/ext4/} source tree for comments
of the form \texttt{TODO:}, \texttt{FIXME:}, and \texttt{WARN\_ON(1)}.
We also read the ATC~2024 paper's future-work section.  We wanted
candidates that (a)~were backed by explicit developer comments in the
code itself, (b)~were not already fixed in newer upstream kernels,
and (c)~could be implemented in a bounded time budget without touching
the on-disk format.

This gave us three promising candidates in order: (1)~a developer TODO
in \texttt{fs/jbd2/commit.c:454} about fast commit blocking too early,
(2)~a developer TODO in \texttt{fs/ext4/mballoc.c:2564} about
synchronous bitmap prefetch completion, and (3)~the xattr fast-commit
gap documented but unpatched in the ATC~2024 paper.

We attempted them in that order.  C1 and C2 produced null measurable
results.  C3 succeeded.  Motivated by C3's success, we then measured
three further candidates against the same criterion (Section~7.1
``characterisation first'') before writing a single line of C4 kernel
code.  Only the fallocate-range path cleared the bar.

The following table summarises all four candidates:

\begin{center}
\small
\begin{tabular}{lcccc}
\toprule
 & \textbf{C1 (barrier)} & \textbf{C2 (bitmap)} & \textbf{C3 (xattr)} & \textbf{C4 (falloc)} \\
\midrule
Slow-path hit rate   & ${\sim}0.1\%$ & ${<}1\%$  & \textbf{100\%}   & \textbf{100\%}  \\
Effect per hit       & ${<}1$~ms     & sub-ms    & 5--11~ms          & 5--7~ms         \\
Measurable on VM?    & No            & No        & Yes ($64{\times}$) & Yes ($62.5{\times}$) \\
Patch size           & 3 lines       & 80 lines  & 108 lines         & 46 lines        \\
\bottomrule
\end{tabular}
\end{center}

C1 and C2 both optimise things that rarely happen in our workloads.
Even if the per-hit savings are real, the aggregate effect drowns in
measurement noise.  C3 optimises every single \texttt{setxattr}; C4
optimises every \texttt{fallocate(COLLAPSE/INSERT\_RANGE)}.  At 100\%
hit rate, even the VM cannot hide a 62--64$\times$ drop in commit count.

% =====================================================================
\section{Candidates 1 and 2: Did Not Measure}
% =====================================================================

Both candidates were implementations of long-standing TODO comments
in the kernel. Both were correct (built, booted, mounted, no
dmesg errors), but neither produced a measurable speedup on our
workloads.

\subsection{C1: fast-commit barrier deferral}

\subsubsection{The developer TODO}

The JBD2 developers themselves left an explicit TODO at
\texttt{fs/jbd2/commit.c} around line~454:

\begin{lstlisting}[language=C]
/*
 * TODO: by blocking fast commits here, we are increasing
 * fsync() latency slightly. Strictly speaking, we don't need
 * to block fast commits until the transaction enters T_FLUSH
 * state. So an optimization is possible where we block new fast
 * commits here and wait for existing ones to complete
 * just before we enter T_FLUSH. That way, the existing fast
 * commits and this full commit can proceed parallely.
 */
\end{lstlisting}

The current code blocks any in-flight fast commit at \texttt{T\_LOCKED}
--- the very first state of a full commit pipeline.  The comment says
we could safely let them run until \texttt{T\_FLUSH}, which comes
several pipeline states later.  If a full commit and a concurrent fast
commit could overlap across \texttt{T\_LOCKED} $\to$ \texttt{T\_SWITCH}
$\to$ \texttt{T\_FLUSH}, those states become effectively free.

\subsubsection{What we implemented}

A three-step change in \texttt{fs/jbd2/commit.c}:

\begin{enumerate}[leftmargin=*]
  \item \textbf{Remove the early drain loop} (lines 444--462). Keep
        only the \texttt{JBD2\_FULL\_COMMIT\_ONGOING} flag set under
        \texttt{j\_state\_lock}, so new fast commits are still blocked
        from \emph{starting}, but in-flight ones may continue.
  \item \textbf{Remove \texttt{journal->j\_fc\_off = 0}} from its
        position at \texttt{T\_LOCKED} (line~472).
  \item \textbf{Insert a new drain + reset} just before
        \texttt{T\_FLUSH} (line~566): wait for in-flight fast commits
        to finish, \emph{then} reset \texttt{j\_fc\_off} to 0.
\end{enumerate}

\paragraph{The hidden \texttt{j\_fc\_off} race we discovered.}
A naive ``just move the drain loop'' fix is \textbf{incorrect}.  The
reason this TODO went unfixed for five-plus years is that the fix has
a subtle hidden constraint: \texttt{j\_fc\_off} is reset to~0 at
\texttt{T\_LOCKED}.  If you move the drain but leave the reset at
\texttt{T\_LOCKED}, an in-flight fast commit's next
\texttt{j\_fc\_off += N} allocation would reuse FC-area blocks that
are supposedly already reserved by that commit.  This silently
corrupts the on-disk fast-commit log.  Our patch moves the reset
alongside the drain so both happen atomically before \texttt{T\_FLUSH}.
This observation --- that you cannot naively move the drain without
also moving the reset --- is a correct insight about the JBD2 fast
commit protocol and is the real technical contribution of C1.

\subsubsection{Correctness analysis}

We manually verified four properties by code inspection (all
file:line references in Linux~6.1.4):

\begin{enumerate}[leftmargin=*]
  \item \textbf{No new fast commits after line~443:}
        \texttt{JBD2\_FULL\_COMMIT\_ONGOING} is set under
        \texttt{j\_state\_lock}; checked by
        \texttt{jbd2\_fc\_begin\_commit} under the same lock
        (\texttt{journal.c:745}).  Mutual exclusion preserved.
  \item \textbf{T\_LOCKED and T\_SWITCH do not touch the FC area:}
        grep confirmed no references to \texttt{j\_fc\_off},
        \texttt{j\_fc\_first}, \texttt{j\_fc\_last}, or
        \texttt{j\_fc\_wbuf} between \texttt{T\_LOCKED} and our new
        drain point.
  \item \textbf{Drain at T\_FLUSH is a sufficient barrier:} after
        the wait loop, \texttt{JBD2\_FAST\_COMMIT\_ONGOING} is clear;
        \texttt{j\_fc\_off = 0} is then safe.
  \item \textbf{Crash recovery is unaffected:} recovery reads
        on-disk data only; it is agnostic to the in-kernel reset
        timing.
\end{enumerate}

\subsubsection{Why it did not work}

The patch is correct and was verified on the VM: it builds, boots,
mounts, and runs without \texttt{dmesg} errors.  The problem was
the measurement.  On our \texttt{fio} \texttt{randwrite --fsync=1
--numjobs=4} workload:

\begin{center}
\begin{tabular}{lrrr}
\toprule
Metric & Stock & Patched & $\Delta$ \\
\midrule
Transactions            & 76   & 75   & --- \\
Avg commit time ($\mu$s) & 6{,}590 & 6{,}661 & $+1.1\%$ (noise) \\
T\_LOCKED phase (ms)   & 0    & 0    & --- \\
T\_FLUSH phase (ms)    & 6    & 6    & --- \\
``0ms locked'' counter  & yes  & yes  & --- \\
\bottomrule
\end{tabular}
\end{center}

The delta is 1\% in the wrong direction, well inside measurement
noise.  The \pth{/proc/fs/jbd2/loopX-8/info} file reported
\texttt{0ms transaction was being locked} on both kernels, meaning
the \texttt{T\_LOCKED} phase already completes in under a
millisecond.  For a fast commit to benefit from the patch, it must
arrive inside that \textless 1~ms window while a full commit is
simultaneously in \texttt{T\_LOCKED}.  The probability of this
overlap is roughly $(1\text{ ms}) / (1000\text{ ms/s}) \times
P(\text{full commit in flight}) \approx 0.1\%$ per \texttt{fsync}.
The optimisation is real; the frequency of the contention in
realistic workloads is too low to show above noise.

An \emph{earlier} ``57\% improvement'' result turned out to be a
setup artifact: the stock benchmark had run for 5~MB while the
patched benchmark ran for 60~s (different system warm-up, different
cache state), and the test filesystem was formatted without
\texttt{-O fast\_commit} (so the drain loop was never reached).
After correcting both, no improvement remained.

We varied \texttt{numjobs} (1, 2, 4, 8 concurrent fsync threads)
and tested all three journaling modes; none closed the gap.  Patch
and postmortem are at \pth{Fast_Commit_Optimization/jbd2-fc-barrier-defer.patch}
and \pth{Fast_Commit_Optimization/CANDIDATE1_postmortem.md}.

\subsection{C2: async bitmap prefetch completion}

\subsubsection{The developer TODO}

In \texttt{fs/ext4/mballoc.c:2564}, the ext4 block allocator
developers left an explicit TODO:

\begin{lstlisting}[language=C]
/*
 * TODO: We should actually kick off the buddy bitmap setup in a work
 * queue when the buffer I/O is completed, so that we don't block
 * waiting for the block allocation bitmap read to finish when
 * ext4_mb_prefetch_fini is called from ext4_mb_regular_allocator().
 */
\end{lstlisting}

\texttt{ext4\_mb\_prefetch\_fini()} runs at the end of every call to
\texttt{ext4\_mb\_regular\_allocator}.  It walks a list of groups for
which bitmaps were prefetched and calls
\texttt{ext4\_mb\_init\_group} on each.  That function blocks on
bitmap I/O if the prefetched read is not yet complete, stalling the
allocator even after it has already found a target block.

\subsubsection{Three implementation shapes considered}

\begin{description}[leftmargin=*]
  \item[A: Simple offload.] Always queue \texttt{ext4\_mb\_init\_group}
        to a workqueue.  Simple, but wastes worker threads blocking
        on I/O when the bitmap read is already done.
  \item[B: True TODO intent.] Hook the bitmap READ bio's
        \texttt{end\_io} callback to schedule the buddy init when I/O
        completes.  Matches the TODO literally; changes the read path
        and raises correctness risk (IRQ context, RCU).
  \item[C: Opportunistic hybrid \textit{(chosen)}.] In
        \texttt{ext4\_mb\_prefetch\_fini}, probe the bitmap buffer: if
        already uptodate, init inline; if not, queue a work item.
        Captures most of (B)'s benefit without touching the read path.
\end{description}

We chose~(C) because it has the lowest correctness risk.  If~(C)
shows measurable benefit, (B) is a natural follow-on; if not, (B)
would not either --- so we save the work.

\subsubsection{Correctness analysis}

Four invariants verified by code inspection:

\begin{enumerate}[leftmargin=*]
  \item \texttt{ext4\_mb\_init\_group} early-returns if
        \texttt{EXT4\_MB\_GRP\_NEED\_INIT} is clear, making races
        between the async worker and synchronous callers harmless ---
        whichever wins, the other no-ops.
  \item The buddy page lock (\texttt{ext4\_mb\_get\_buddy\_page\_lock})
        serialises the actual cache work between any two callers.
  \item \texttt{flush\_workqueue} + \texttt{destroy\_workqueue} in
        \texttt{ext4\_mb\_release()} drain all pending workers before
        group metadata is freed on unmount.
  \item On init failure, the error path destroys both the slab and
        the workqueue we created; \texttt{ext4\_mb\_release} is not
        called for failed mounts, so there is no double-destroy.
\end{enumerate}

\subsubsection{Bugs caught during code review}

\paragraph{Blocker 1 --- \texttt{static noinline\_for\_stack} orphaning.}
Our first version inserted the new \texttt{struct
ext4\_bitmap\_init\_work} and helper functions \emph{between}
\texttt{static noinline\_for\_stack} (line~1363) and
\texttt{int ext4\_mb\_init\_group(...)} (line~1364).  C grammar
binds \texttt{static noinline\_for\_stack} to the \emph{next}
declaration, which became our struct, not the function.
Consequences: \texttt{ext4\_mb\_init\_group} silently lost
\texttt{static} linkage (symbol namespace pollution) and its
\texttt{noinline\_for\_stack} annotation (removing the stack
isolation needed on deeply-nested allocator call chains with KASAN).
\textbf{Fix:} moved the inserted block to before the ``Locking note''
comment, restoring adjacency with the function definition.

\paragraph{Blocker 2 --- resource leak on OOM in \texttt{ext4\_mb\_init}.}
We originally placed the slab and workqueue creation \emph{after}
\texttt{ext4\_mb\_init\_backend()} succeeded.  If our
\texttt{alloc\_workqueue()} call then failed, the error path jumped
to \texttt{out\_free\_locality\_groups}, which could not undo
the backend's work, leaking \texttt{s\_buddy\_cache} and
\texttt{s\_group\_info} on that specific OOM path.
\textbf{Fix:} moved creation to \emph{before}
\texttt{ext4\_mb\_init\_backend}.  If creation fails the backend has
not run yet --- nothing to clean up.

\paragraph{Build-time issue --- forward declaration.}
After moving the helper functions ahead of
\texttt{ext4\_mb\_init\_group}, the build failed: our
\texttt{ext4\_mb\_async\_init\_worker} calls
\texttt{ext4\_mb\_init\_group()} before the compiler sees its
\texttt{static} definition.  \textbf{Fix:} one-line forward
declaration above the new block.  Neither review stage caught this
because both reviewers read the patched source textually, not with
compiler eyes.  Build-and-boot verification is non-negotiable.

\subsubsection{Why it did not work}

The patch is correct and the kernel boots without errors.  The
problem was again measurement.

\paragraph{Fallocate microbench --- no signal.}
We ran a fresh 16~GB loop device, \texttt{mkfs.ext4 -F -O
fast\_commit}, and iterated 32$\times$ \texttt{fallocate -l 256M}
with sync between.  Three full repeats (96 samples):

\begin{center}
\begin{tabular}{lrrr}
\toprule
Metric & Stock & Patched & $\Delta$ \\
\midrule
Mean ($\mu$s)   & 8{,}629 & 9{,}006 & $+4.4\%$ worse \\
Median ($\mu$s) & 8{,}422 & 8{,}703 & $+3.3\%$ worse \\
p95 ($\mu$s)    & 10{,}682 & 12{,}426 & $+16\%$ worse \\
p99 ($\mu$s)    & 14{,}998 & 13{,}647 & $-9\%$ better \\
Stdev ($\mu$s)  & 1{,}051  & 1{,}180  & $+12\%$ worse \\
\bottomrule
\end{tabular}
\end{center}

All deltas are within 1$\sigma$: a null result.  The reason: a
16~GB filesystem with default 128~MB groups has only 128 groups.
After the first few 256~MB allocations the bitmaps are warm;
\texttt{EXT4\_MB\_GRP\_NEED\_INIT} is cleared; subsequent
allocations reuse cached bitmaps and never reach the
\texttt{prefetch\_fini} async branch.  In our workload the cold-init
path is hit on less than 1\% of allocations.

\paragraph{VM variance discovery.}
An \texttt{fio} throughput benchmark (8~GB write, 4 jobs) showed a
promising $+10.5\%$ on the first composite run.  A focused 5-repeat
test revealed the truth:

\begin{center}
\begin{tabular}{lr}
\toprule
Run & BW (MiB/s) \\
\midrule
1 & 774 \\
2 & 655 \\
3 & 203 \\
4 & 203 \\
5 & 190 \\
\bottomrule
\end{tabular}
\end{center}

Mean~$405 \pm 286$~MiB/s --- a \textbf{70\% coefficient of
variation} within a single kernel.  The VM's loop device passes
writes through the host page cache; the first two runs are
memory-bound, later runs actually hit host disk.  At 70\% CV no
patch smaller than a $2\times$ speedup can be measured.  The
initial $+10.5\%$ was an order-of-measurement artifact.

Patch and postmortem are at \pth{Fast_Commit_Optimization/mballoc-async-prefetch.patch}
and \pth{Fast_Commit_Optimization/CANDIDATE2_postmortem.md}.

\paragraph{Lesson.} A patch is only as useful as the fraction of
operations that exercise the path it changes.  After C1 and C2 we
adopted a hard rule: \emph{only ship a patch if a microbench can
show that the slow path it removes is hit on a substantial fraction
of the target operation, and the measurement noise floor is lower
than the expected effect.}  C3 and C4 were both validated against
this rule before any kernel code was written.

% =====================================================================
\section{Candidate 3: Fast Commit Support for Inline xattrs}
% =====================================================================

\subsection{Intuition}

Every \texttt{setxattr} or \texttt{removexattr} call on a
fast-commit-enabled ext4 filesystem forces a full JBD2 commit.
The cost is paid on every SELinux relabel, every POSIX ACL change,
and every short \texttt{user.*} attribute write, which makes the
slow path effectively 100\% hit on xattr-touching workloads. The
ATC~2024 paper~\cite{atc24-fc} flags this gap but provides no
patch.

This matters because extended attributes are used pervasively.
SELinux labels every file via \texttt{security.selinux}.  POSIX
ACLs use \texttt{system.posix\_acl\_access}.  Package managers
write \texttt{user.*} attributes for checksums.
\texttt{systemd-tmpfiles} writes xattrs during boot.  On any
fast-commit filesystem, every one of these operations silently
reverts to the expensive full-commit path even though the fast
commit area sits idle.

\subsection{Where xattrs live in ext4}

ext4 stores xattrs in three locations:

\begin{enumerate}[leftmargin=*]
  \item \textbf{Inline} --- in the inode body itself, between
        \texttt{EXT4\_GOOD\_OLD\_INODE\_SIZE + i\_extra\_isize}
        and the end of the inode.  On a 256-byte inode this gives
        roughly 80~bytes for xattrs.  Most SELinux labels, POSIX
        ACLs, and short \texttt{user.*} xattrs fit here.
  \item \textbf{Block} --- in a separate 4~KB block pointed to by
        \texttt{i\_file\_acl}.  Used when xattrs do not fit inline.
  \item \textbf{EA inode} --- each xattr gets its own small inode.
        Used when the value is very large and the \texttt{ea\_inode}
        feature is enabled.
\end{enumerate}

Our patch handles case~(1).  For cases~(2) and~(3) we keep the
full-commit fallback.  This covers nearly all real-world xattr
operations because SELinux and ACLs are always short and fit inline.

\subsection{Code analysis}

The fallback is fired by an unconditional call at the end of
\texttt{ext4\_xattr\_set\_handle} in \texttt{fs/ext4/xattr.c}
(line~2422):

\begin{lstlisting}[language=C]
ext4_fc_mark_ineligible(inode->i_sb, EXT4_FC_REASON_XATTR, handle);
\end{lstlisting}

That single line kills fast commit for the entire transaction,
regardless of whether the xattr was inline, block, or EA-inode.

A second \texttt{mark\_ineligible} call sits at line~2500 inside
the wrapper \texttt{ext4\_xattr\_set}, \emph{after}
\texttt{ext4\_journal\_stop}, with \texttt{handle==NULL}.  With a
null handle this call marks the current running JBD2 transaction
ineligible, suppressing fast commit on any later \texttt{fsync}
sharing that transaction.  This second call is more subtle ---
it is in a different function from the one our patch changes, and
finding it required tracing every call site of
\texttt{ext4\_fc\_mark\_ineligible} in the xattr code.  Missing
it would have made the patch a no-op for end users even though the
change inside \texttt{ext4\_xattr\_set\_handle} was correct.
We deleted it after confirming that our new conditional at
line~2422 already handles every case on which the inode is actually
modified.

\subsection{Design and implementation}

Two changes, both small.

\begin{enumerate}[leftmargin=*]
  \item \textbf{Expand the FC inode log.}
        \texttt{ext4\_fc\_write\_inode} in
        \texttt{fs/ext4/fast\_commit.c:863} currently logs the
        inode's raw bytes up to
        \texttt{EXT4\_GOOD\_OLD\_INODE\_SIZE + i\_extra\_isize}.
        That stops just before the inline xattr region.  We add a
        two-line conditional that expands the log to
        \texttt{EXT4\_INODE\_SIZE} when the inode has inline
        xattrs, detected via the pre-existing
        \texttt{EXT4\_STATE\_XATTR} flag that ext4 already sets
        and clears correctly in its xattr code.

        Replay needs no change: \texttt{ext4\_fc\_replay\_inode}
        already \texttt{memcpy}s \texttt{fc\_len} bytes from the
        log into the inode, and the validator already accepts
        records up to \texttt{s\_inode\_size}.  When the patched
        kernel logs a larger inode, the same replay code restores
        it correctly.  This is the main reason the patch came out
        so small --- we reused the entire existing replay path.

  \item \textbf{Skip the ineligibility call when safe.} In
        \texttt{ext4\_xattr\_set\_handle}, add a local
        \texttt{bool touched\_block = false} at the top. Set it
        to true on every path that calls
        \texttt{ext4\_xattr\_block\_set} or takes the EA-inode
        path (\texttt{i.in\_inode = 1}).  Replace the unconditional
        \texttt{mark\_ineligible} at line~2422 with:
\begin{lstlisting}[language=C]
if (error || touched_block || !ext4_has_feature_xattr(inode->i_sb))
    ext4_fc_mark_ineligible(inode->i_sb,
                            EXT4_FC_REASON_XATTR, handle);
\end{lstlisting}
        and delete the wrapper's redundant call at line~2500.

        The three conditions each have a specific role:
        \begin{itemize}[leftmargin=*]
          \item \texttt{error}: the operation failed; fall back to be safe.
          \item \texttt{touched\_block}: we modified a separate xattr
                block or an EA inode; fast commit does not log these,
                so a full commit is required.
          \item \texttt{!ext4\_has\_feature\_xattr(sb)}: this is the
                \emph{very first} xattr on a freshly-mkfs'd filesystem.
                The superblock xattr feature bit has just been set in
                memory but not yet persisted.  If we fast-committed now
                and the machine crashed, the replayer would not see the
                bit and would not know to interpret the inode's inline
                region as xattr data.  The first xattr is therefore
                always a full commit by design.
        \end{itemize}
\end{enumerate}

The full patch is 108 lines (\pth{Fast_Commit_Optimization/fc-inline-xattr.patch}):
\texttt{+9/-0} in \texttt{fast\_commit.c} (the new conditional with
its comment), and \texttt{+27/-4} in \texttt{xattr.c} (the
\texttt{touched\_block} declaration, five assignments to it on
each block-touching path, the new conditional at line~2422, the
removal of the redundant call at line~2500, and explanatory
comments at both sites).

\subsection{Correctness}

\begin{itemize}[leftmargin=*]
  \item Three crash-recovery scripts
        (\texttt{c3\_crash\_test\_\{a,b,c\}.sh}). Test~A: 100
        inline xattrs, lazy unmount, remount, count
        (\textbf{PASS}, 100/100). Test~B: 100 set, 50 even
        removed, lazy unmount, remount, verify exactly the 50
        odd ones remain (\textbf{PASS}). Test~C: a 668-byte xattr
        value (well over the inline limit so the xattr block path
        is used), lazy unmount, remount, full value recovered
        (\textbf{PASS}, confirming the fallback branch still runs
        the full commit when needed).
  \item xfstests subset
        \texttt{generic/\{062,118,300,337,454,455,473,482\}} +
        \texttt{ext4/032}. Of the nine tests, three skip due to
        missing features (reflink for \texttt{generic/118},
        \texttt{LOGWRITES\_DEV} for \texttt{generic/455} and
        \texttt{482}). Of the six that ran, five pass and
        \texttt{generic/473} fails identically on stock and
        patched (a pre-existing 6.1.4 issue, not a regression).
        Logs at \pth{Fast_Commit_Optimization/xfstests_c3/}.
\end{itemize}

\subsection{Results}

\paragraph{VM (VirtualBox, 11th Gen Intel Core i7-11800H, 4 cores,
5.7 GiB RAM; 256 MB loop fs with \texttt{-O fast\_commit}; 5000
\texttt{fsetxattr}+\texttt{fsync} on a single inline-sized value).}

\begin{center}
\begin{tabular}{lrrr}
\toprule
Metric & Stock & Patched & Change \\
\midrule
Wall time (5000 ops) & 31{,}102~ms & 22{,}734~ms & $-26.9\%$ \\
Per-op latency       & 6{,}220~$\mu$s & 4{,}546~$\mu$s & $-26.9\%$ \\
JBD2 full commits    & 5{,}000 & \textbf{78} & $\mathbf{-98.4\%}$ \\
Commits per op       & 1.000   & 0.016       & --- \\
\bottomrule
\end{tabular}
\end{center}

The 64$\times$ reduction in full commits (5000~$\to$~78) is the
clean architectural signal.  Our patch engages the fast commit
path on 98.4\% of operations, exactly as designed.  The remaining
1.6\% are expected: one full commit for the warmup xattr
(feature-bit persistence, by design), and a handful from JBD2's
natural 5-second timer and fast-commit-area wraparounds.

The 27\% wall-time reduction is smaller than the commit-count
drop might suggest.  On the VM each \texttt{fsync} takes about
5~ms end-to-end, of which roughly 1.7~ms is the JBD2 full commit
we eliminate; the remaining~3.3~ms is VFS, syscall entry/exit,
and the host-backed loop device's write-behind path --- none of
which our patch touches.  On real hardware this ratio shifts.

\paragraph{Bare-metal (11th Gen Intel Core i3-1115G4, 4 cores,
7.5 GiB RAM, real SATA SSD; 3 repeats).}

\begin{center}
\begin{tabular}{lrrr}
\toprule
Metric & Stock & Patched & Change \\
\midrule
Wall time (mean)     & 57{,}627~ms & 30{,}089~ms & $-47.8\%$ \\
Stdev                & 3{,}225~ms  & 2{,}088~ms  & (${\sim}6\%$ CV) \\
Per-op latency       & 11{,}525~$\mu$s & 6{,}017~$\mu$s & $-47.8\%$ \\
JBD2 full commits    & 5{,}000 & \textbf{78} & $\mathbf{-98.4\%}$ \\
\bottomrule
\end{tabular}
\end{center}

Two observations.  First, the transaction count reduction is
\textbf{byte-identical} to the VM result: 64$\times$,
5000~$\to$~78.  This is the architectural claim of the patch;
it reproduces across machines with different CPUs, memory, and
storage.  Second, the wall-time speedup is nearly doubled on the
bare-metal machine: 48\% versus 27\%.  On a real SATA SSD the
VFS and syscall overhead form a smaller share of per-op cost, so
eliminating the JBD2 commit portion yields a larger relative gain.
The 3-repeat coefficient of variation is about 6\%, a usable
measurement setup; the result is not a one-off artifact.

\paragraph{Non-xattr regression check.}
To confirm the patch does not regress unrelated workloads we ran
\texttt{fio --rw=randwrite --bs=4k --fsync=1 --numjobs=4} for 30~s
on the patched kernel.  Stock: 530~IOPS.  Patched: 528~IOPS.
Delta 0.4\%, well inside noise.  No regression.

% =====================================================================
\section{Candidate 4: Fast Commit Support for fallocate Range Operations}
% =====================================================================

\subsection{Intuition}

After C3 succeeded we asked a simple question: is there another
slow path in ext4 with the same shape?  That is, an operation that
(a)~is used regularly on a fast-commit filesystem, (b)~hits
\texttt{ext4\_fc\_mark\_ineligible} on 100\% of calls, and (c)~has
a simple enough code path that the fix is a small patch, not a
rewrite.

We found one: the \texttt{FALLOC\_FL\_COLLAPSE\_RANGE} and
\texttt{FALLOC\_FL\_INSERT\_RANGE} modes of \texttt{fallocate(2)}.
Collapse removes a middle slice of a file and shifts following data
left; insert adds a zero-filled hole at an offset and shifts
following data right.  Both are used by log rotators, database log
compactors, and VM disk image tools.  Both call
\texttt{ext4\_fc\_mark\_ineligible(EXT4\_FC\_REASON\_FALLOC\_RANGE)}
unconditionally and force a full JBD2 commit on every invocation.

\subsection{Characterisation first}

The cost of our first two failures was implementing patches before
confirming the slow path was hit hard enough to measure.  For C4
we flipped the order: measure first, write kernel code second.

We ran three candidate microbenchmarks on the C3 kernel to compare
three possible targets: fallocate COLLAPSE/INSERT range; large
(4~KB) \texttt{setxattr} values that spill to a separate xattr
block (the case C3 does not cover); and many-thread concurrent
\texttt{fsync} (the CJFS-style compound flush target).  Three runs
each, $N=1000$ ops (or 100 per thread at $T=16$).

The selection rule: only implement if the microbench shows
\texttt{tx/op $\ge$ 0.5} --- at least one full commit per op.
Anything below means the slow path is already being coalesced.

\begin{center}
\begin{tabular}{lrrl}
\toprule
Workload & ms/op & tx/op & Signal \\
\midrule
fallocate COLLAPSE\_RANGE & 7.03 & 1.001 & \textbf{strong} \\
fallocate INSERT\_RANGE   & 7.01 & 1.001 & \textbf{strong} \\
xattr block (4~KB value)  & 7.23 & 1.001 & strong but narrow \\
concurrent fsync $T=16$   & 1.21 & 0.003 & weak \\
\bottomrule
\end{tabular}
\end{center}

Fallocate COLLAPSE/INSERT showed the same per-op signal as C3's
xattr case: every single op triggers exactly one full commit.
The xattr-block path showed equal per-op cost but covers a
minority of \texttt{setxattr} calls in practice --- typical
SELinux and ACL values are short and already handled by C3.
Concurrent fsync was already well-coalesced by JBD2's group commit:
the tx/op ratio drops sharply as threads increase (0.020 at $T=1$
down to 0.003 at $T=16$), so a CJFS-style port would not help much
on 6.1.4.  We chose fallocate for C4 because it is a clean,
independent code path rather than an extension of C3 --- the patch
adds a second, separate architectural reduction rather than a
refinement of the first.

\subsection{Code analysis}

Two call sites in \texttt{fs/ext4/extents.c} mark fallocate
range operations as ineligible, both right after
\texttt{ext4\_journal\_start} and before the actual extent-tree
work:

\begin{lstlisting}[language=C]
/* line 5363, inside ext4_collapse_range */
ext4_fc_mark_ineligible(sb, EXT4_FC_REASON_FALLOC_RANGE, handle);

/* line 5508, inside ext4_insert_range */
ext4_fc_mark_ineligible(sb, EXT4_FC_REASON_FALLOC_RANGE, handle);
\end{lstlisting}

The rest of each function modifies the extent tree (removes a
range, inserts a hole, shifts extents left or right) and updates
\texttt{i\_size} and the timestamps.

Crucially, ext4 \emph{already} has fast-commit tags for the
underlying extent-level changes: \texttt{FC\_TAG\_ADD\_RANGE}
and \texttt{FC\_TAG\_DEL\_RANGE}, emitted by
\texttt{ext4\_fc\_write\_inode\_data} in \texttt{fast\_commit.c}.
The per-inode \texttt{FC\_TAG\_INODE} tag already carries the
updated \texttt{i\_size}.  All the information needed to reproduce
a collapse or insert on the replay side is already representable
in the existing tag vocabulary.

The \texttt{FALLOC\_RANGE} fallback exists for one reason only:
the collapse/insert paths never call \texttt{ext4\_fc\_track\_range}
on the logical blocks they modify.  Without that call the
fast-commit commit walker does not know to include the shifted
region when it walks the extent tree at commit time.

\subsection{Design and implementation}

Replace each of the two \texttt{mark\_ineligible} calls with a
call to \texttt{ext4\_fc\_track\_range} over the affected
logical-block range, computed at the point where the journal
handle is open (so \texttt{inode->i\_size} still holds the
pre-modification size):

\begin{lstlisting}[language=C]
/* collapse path: track from punch_start to current EOF */
ext4_fc_track_range(handle, inode, punch_start,
    ((inode->i_size + inode->i_sb->s_blocksize - 1) >>
     EXT4_BLOCK_SIZE_BITS(inode->i_sb)) - 1);

/* insert path: track from offset_lblk to post-grow EOF */
ext4_fc_track_range(handle, inode, offset_lblk,
    ((inode->i_size + len + inode->i_sb->s_blocksize - 1) >>
     EXT4_BLOCK_SIZE_BITS(inode->i_sb)) - 1);
\end{lstlisting}

For collapse, the affected range runs from \texttt{punch\_start}
up to the current EOF.  For insert, we add \texttt{len} to cover
the grown region.  At fsync time, the standard fast-commit commit
runs \texttt{ext4\_fc\_write\_inode\_data}, which walks the
post-shift extent tree over the tracked range, emits
\texttt{FC\_TAG\_ADD\_RANGE} for each shifted extent at its new
logical position, and \texttt{FC\_TAG\_DEL\_RANGE} for any
resulting hole (new in the insert case; at the old tail in the
collapse case).  \texttt{FC\_TAG\_INODE} is emitted by the
\texttt{ext4\_mark\_inode\_dirty} call that already exists later in
each function and carries the new \texttt{i\_size}.

\paragraph{Replay correctness.} The non-obvious part is that
\texttt{ext4\_fc\_replay\_del\_range} frees the physical blocks of
the range it removes via \texttt{ext4\_mb\_mark\_bb(\ldots,~0)}.
In an insert operation, only the logical positions move --- the
physical data blocks stay put.  \texttt{DEL\_RANGE} replay
therefore transiently marks those blocks free, and
\texttt{ADD\_RANGE} replay re-binds them at their new logical
position.  The intermediate bitmap state is wrong; the final state
is reconciled by the pre-existing
\texttt{ext4\_fc\_set\_bitmaps\_and\_counters} pass, which runs at
the end of every replay and re-walks all modified inodes' extent
trees to re-mark currently-reachable blocks as used.  This pass is
what makes the design safe without requiring new tag types or
changes to replay logic.

The full patch is 46 lines (\pth{Fast_Commit_Optimization/fc-fallocate-range.patch}),
two hunks in \texttt{fs/ext4/extents.c}.  No change to
\texttt{fast\_commit.c}, no new tag types, no on-disk format change,
no new compat feature bit.  Same philosophy as C3: reuse existing
primitives rather than add new ones.

\subsection{Correctness}

\paragraph{Compilation.}  \texttt{fs/ext4/extents.o} builds clean
at kernel warning level~2.  The full kernel build with
\texttt{LOCALVERSION=-cs614-c4-patched} completes with zero errors.

\paragraph{Smoke test.}  On the booted C4 kernel we created a file
with distinct per-block byte patterns, performed a collapse
followed by an insert, read back the shifted region, and compared
md5 checksums.  The output matched the expected post-operation
state computed in Python.

\paragraph{Custom crash-recovery tests.}  Three tests at
\pth{Fast_Commit_Optimization/c4_crash_test_{a,b,c}.sh}:
\begin{itemize}[leftmargin=*]
  \item Test~A: write a 1~MB file with distinct per-block patterns,
        \texttt{fsync}, \texttt{fallocate -c} to collapse 32 blocks,
        \texttt{fsync}, simulate crash via lazy unmount, remount,
        compare md5 and size against expected post-collapse state.
        \textbf{PASS}.
  \item Test~B: similar, but \texttt{fallocate -i} to insert 16
        blocks at offset~16.  Expected state: unchanged head,
        zero-filled hole in the inserted region, original tail
        shifted right.  \textbf{PASS}.
  \item Test~C: three interleaved operations (collapse, insert,
        collapse) each followed by \texttt{fsync}, then crash.
        Expected md5 computed in pure Python from the same sequence.
        \textbf{PASS}.  This is the strongest test because it
        accumulates multiple \texttt{ADD\_RANGE} and
        \texttt{DEL\_RANGE} tags for the same inode in the FC log,
        stressing the multi-op replay path.
\end{itemize}

\paragraph{xfstests subset.}  Same 9-test subset as C3
(\texttt{ext4/032} and \texttt{generic/\{062, 118, 300, 337, 454,
455, 473, 482\}}).  Six tests ran on our setup.  \textbf{5 pass}.
The one failure is \texttt{generic/473} with a diff byte-identical
to the failure on the C3 kernel and on stock 6.1.4 --- a
pre-existing ext4 bug, not a C4 regression.  Critically,
\texttt{generic/062, 337, 454} (xattr-heavy tests) all pass,
confirming C4 does not regress C3.

\subsection{Results}

\paragraph{VM (VirtualBox i7-11800H; 256~MB loop ext4 with
\texttt{-O fast\_commit}; 1000 ops + per-op \texttt{fsync};
3 runs per mode).}  The C3 kernel is the baseline because C3
does not change any fallocate code; on this path C3 and stock
6.1.4 are architecturally identical.

\begin{center}
\begin{tabular}{lrrr}
\toprule
Metric & Baseline (C3) & Patched (C4) & Change \\
\midrule
\multicolumn{4}{l}{\textit{COLLAPSE\_RANGE}} \\
Wall time (mean)     & 7{,}029~ms       & 5{,}423~ms       & $-22.8\%$ \\
Per-op latency       & 7{,}029~$\mu$s   & 5{,}423~$\mu$s   & $-22.8\%$ \\
JBD2 full commits    & 1{,}001 & \textbf{16} & $\mathbf{-98.4\%}$ \\
Commits per op       & 1.001   & 0.016       & --- \\
\midrule
\multicolumn{4}{l}{\textit{INSERT\_RANGE}} \\
Wall time (mean)     & 7{,}006~ms       & 5{,}189~ms       & $-25.9\%$ \\
Per-op latency       & 7{,}006~$\mu$s   & 5{,}189~$\mu$s   & $-25.9\%$ \\
JBD2 full commits    & 1{,}001 & \textbf{16} & $\mathbf{-98.4\%}$ \\
Commits per op       & 1.001   & 0.016       & --- \\
\bottomrule
\end{tabular}
\end{center}

Both modes show a 62.5$\times$ reduction in JBD2 full commits
(1001~$\to$~16 out of 1000~ops).  The remaining 16 commits are
the expected cases: JBD2's 5-second timer, warmup flushes, and
occasional FC-area wraparounds.  The wall-time gain is about 24\%,
smaller than the commit-count drop suggests because fallocate range
operations have a significant non-journal cost: extent-tree
rebalance in \texttt{ext4\_ext\_remove\_space} and
\texttt{ext4\_ext\_shift\_extents}.  Our patch does not touch that
work.

\paragraph{Bare-metal (i3-1115G4, 4 cores, 7.5~GiB RAM, real
SATA SSD; 3 repeats; baseline kernel \texttt{6.1.4-cs614-c3-patched},
patched kernel \texttt{6.1.4-cs614-c4-patched} with both C3 and
C4 applied).}

\begin{center}
\begin{tabular}{lrrr}
\toprule
Metric & Baseline (C3) & Patched (C3+C4) & Change \\
\midrule
\multicolumn{4}{l}{\textit{COLLAPSE\_RANGE}} \\
Wall time (mean)     & 12{,}445~ms     & 7{,}032~ms      & $-43.5\%$ \\
Stdev                & 970~ms          & 39~ms            & (${\sim}0.5\%$ CV) \\
Per-op latency       & 12{,}445~$\mu$s & 7{,}032~$\mu$s  & $-43.5\%$ \\
JBD2 full commits    & 1{,}001 & \textbf{16} & $\mathbf{-98.4\%}$ \\
\midrule
\multicolumn{4}{l}{\textit{INSERT\_RANGE}} \\
Wall time (mean)     & 11{,}539~ms     & 6{,}680~ms      & $-42.1\%$ \\
Stdev                & 541~ms          & 142~ms           & (${\sim}2\%$ CV) \\
Per-op latency       & 11{,}539~$\mu$s & 6{,}680~$\mu$s  & $-42.1\%$ \\
JBD2 full commits    & 1{,}001 & \textbf{16} & $\mathbf{-98.4\%}$ \\
\bottomrule
\end{tabular}
\end{center}

Two observations match the C3 pattern exactly.  First, the
transaction-count reduction is \textbf{byte-identical} to the VM:
1001~$\to$~16 in both modes.  This is the architectural claim of
the patch; it reproduces across machines without any tuning.
Second, the wall-time speedup is \textbf{nearly doubled} on
bare-metal: 43--44\% vs 23--26\% on the VM.  The reason is the
same as for C3: on a real SSD the non-journal overhead (VFS,
extent-tree manipulation) forms a smaller fraction of total op
time, so eliminating the JBD2 full commit yields a proportionally
larger gain.  Variance on the patched runs is exceptionally low:
39~ms stdev on a 7000~ms mean (0.5\% CV) for COLLAPSE,
142~ms on 6680~ms (2\% CV) for INSERT.  The result is not a
one-off artifact.

\paragraph{Non-fallocate regression check.}
Our xfstests subset already catches regressions on the most-travelled
metadata paths, and all relevant tests pass.  As an additional
check we verified the C4 kernel still shows 78 full commits on
the 5000-iteration \texttt{setxattr} loop from C3.  C3's benefit
is fully preserved when C4 is stacked on top.

% =====================================================================
\section{Discussion: Why C3 and C4 Worked Where C1 and C2 Did Not}
% =====================================================================

The difference comes down to one metric: the hit rate of the slow
path being optimised.

\begin{center}
\small
\begin{tabular}{lcccc}
\toprule
 & \textbf{C1 (barrier)} & \textbf{C2 (bitmap)} & \textbf{C3 (xattr FC)} & \textbf{C4 (falloc FC)} \\
\midrule
Slow-path hit rate   & ${\sim}0.1\%$ & ${<}1\%$  & \textbf{100\%}   & \textbf{100\%}  \\
Effect per hit       & ${<}1$~ms     & sub-ms    & 5--11~ms          & 5--7~ms         \\
Measurable on VM?    & No            & No        & Yes ($64{\times}$) & Yes ($62.5{\times}$) \\
Patch lines          & 3             & 80        & 108               & 46              \\
\bottomrule
\end{tabular}
\end{center}

C1 and C2 both optimise events that almost never happen in our
workloads.  Even if the per-hit savings are real, the aggregate
effect is buried in noise.  A 1~ms optimisation that triggers
0.1\% of the time contributes 1~$\mu$s of average improvement ---
not a measurable result on any VM.

C3 and C4 optimise events that happen on \emph{every single
invocation} of the target operation.  At a 100\% hit rate, even
our noisy VM --- where the fio throughput benchmark had a 70\% CV
on a single kernel --- cannot hide a 62--64$\times$ drop in
transaction count.

The two patches together remove two of the ten fast-commit
fallback reasons documented in \texttt{fs/ext4/fast\_commit.h}:
\texttt{EXT4\_FC\_REASON\_XATTR} (C3, inline case) and
\texttt{EXT4\_FC\_REASON\_FALLOC\_RANGE} (C4, collapse and insert
modes).  They are independent patches that touch separate code
paths (\texttt{xattr.c} and \texttt{extents.c}); they compose
without conflict and their benefits stack.

\paragraph{Why the wall-time gain is smaller than the commit-count gain.}
In both patches the commit-count reduction (${\sim}64\times$ for
C3, $62.5\times$ for C4) is much larger than the wall-time gain
(${\sim}1.4\times$ on the VM, ${\sim}2\times$ on bare-metal).  The
gap exists because \texttt{fsync} has fixed overhead we do not
touch: VFS, syscall, device write-behind, and (for C4) extent-tree
manipulation in \texttt{ext4\_ext\_remove\_space} and
\texttt{ext4\_ext\_shift\_extents}.  Even a perfect journal
optimisation that made commits free would not eliminate this floor.
The gap closes on real hardware because the non-journal share of
each \texttt{fsync} is smaller there --- exactly what we observed:
C3 went from 27\% on the VM to 48\% on bare-metal; C4 went from
23--26\% to 42--44\%.  The two-step factor between platforms is
consistent across both patches.

\paragraph{Scope limits and future work.}
Still in the full-commit fallback after C3 and C4:
\begin{itemize}[leftmargin=*]
  \item \textbf{Block xattrs and EA-inodes}: would require new tag
        types (\texttt{FC\_TAG\_XATTR\_SET},
        \texttt{FC\_TAG\_XATTR\_DEL}) and e2fsprogs changes to
        match the byte-identical format constraint in
        \texttt{fast\_commit.h}.  Our characterisation confirmed a
        strong per-op signal (7.23~ms, tx/op~=~1.001), but the
        real-world hit rate is much smaller since typical SELinux
        and ACL values are short.
  \item \textbf{CROSS\_RENAME, RENAME\_DIR}: replay cannot yet
        safely update \texttt{..} dirents.
  \item \textbf{FALLOC\_FL\_PUNCH\_HOLE, ordinary fallocate}: already
        handled upstream; not a gap.
\end{itemize}

% =====================================================================
\section{Candidate 5: Lockless CAS for jbd2\_log\_wait\_commit}
% =====================================================================


\subsection{Analyzing the lock latency of JBD2}

High-concurrency applications rely on the \texttt{fsync()} system
call to safely flush in-memory data to the physical disk. In the
ext4 filesystem this is handled by the JBD2 (Journaling Block
Device 2) layer. During a transaction, any application threads
that request a synchronous flush are put to sleep on the
\texttt{j\_wait\_done\_commit} queue. A separate, single
background daemon thread is responsible for actually writing the
data to the disk.

During our analysis using benchmarks that support multhithreading capabilty, we identified a major
scaling bottleneck within the \texttt{jbd2\_log\_wait\_commit()}
function. To manage the sleeping threads safely, the native kernel
places a global read-write lock (\texttt{j\_state\_lock}) around
the code that checks if the commit is finished.

When the background daemon finishes writing to the disk, it sends
a wakeup signal to all the sleeping application threads at the
exact same time. The problem occurs immediately after they wake
up: every single thread must acquire the exact same
\texttt{j\_state\_lock} to verify if their specific data was
committed. Because only a limited number of operations can modify
the lock code at once, the threads spend a massive amount of CPU
time simply waiting in a queue to acquire the lock, severely
delaying their return to the user application.

\begin{lstlisting}[caption={The Native ext4 Implementation (Problematic)}]
int jbd2_log_wait_commit(journal_t *journal, tid_t tid)
{
    int err = 0;

    /* Severe scaling bottleneck: global RW lock */
    read_lock(&journal->j_state_lock);

    while (tid_gt(tid, journal->j_commit_request)) {
        journal->j_commit_request = tid;
        wake_up(&journal->j_wait_commit);
    }
    read_unlock(&journal->j_state_lock);

    wait_event(journal->j_wait_done_commit,
           !tid_gt(tid, journal->j_commit_sequence));

    return err;
}
\end{lstlisting}

\subsection{Identifying lock contention time}

To analyze this wait-time overhead without artificially slowing down the system using heavy sampling tools, we built a custom kernel module. Rather than modifying the filesystem directly, this module passively hooks into the entry and return bounds to track the total time of execution of both the \texttt{fsync()} system paths and the \texttt{jbd2\_log\_wait\_commit()} function.

Inside the module, tracking variables safely keep count of exactly how many total nanoseconds the computer spends running a full filesystem flush. At the exact same time, they strictly measure how many of those nanoseconds were wasted simply waiting in line to acquire the \texttt{j\_state\_lock}. When you remove the module utilizing the \texttt{rmmod} command, it prints a clear statistical summary table directly to the kernel's \texttt{dmesg} log.

By analyzing the resulting kernel \texttt{dmesg} logs across 4 fully independent test runs, the output consistently mapped the latency bottleneck. Reading the custom latency tracker empirically proved that simply waiting to acquire the \texttt{j\_state\_lock} was severely restricting the filesystem's ability to process parallel system operations:

\begin{lstlisting}[caption={Sample \texttt{dmesg} kernel ring buffer output from the latency tracker on \texttt{rmmod}.},language=bash]
============================================
FSYNC AND WAIT-COMMIT LATENCY  (device: nvme0n1p6)
Mode: BASELINE
============================================
PHASE                       COUNT TOTAL_ms AVG_us MIN_us MAX_us
ext4_sync_file              2040  20021    9814   4640   17452
file_write_and_wait_range   2040  4476     2194   75     8446
jbd2_log_wait_commit        2040  15536    7615   4299   14780
jbd2_commit_transaction     509   2531     4974   3310   12367
SUMMARY
jbd2_trace: 1. MAX total execution time: 2539 ms
jbd2_trace: 2. Total aggregated fsync CPU time: 20021 ms
jbd2_trace: 3. Total aggregated time taken on locks: 115 ms
jbd2_trace: 4. Total count of fsyncs: 2040
jbd2_trace: done.
\end{lstlisting}

As demonstrated by the data, while the overall execution path (\texttt{ext4\_sync\_file}) required 20,021 aggregate milliseconds of CPU time to strictly complete the 2,040 file sync operations, the isolated \texttt{jbd2\_log\_wait\_commit} phase independently consumed 15,536 CPU milliseconds of that total limit. So this gives us a direction, if we can try to optimize the jbd2\_log\_wait\_commit function we can have performance boost in journalling. 

\subsection{Mitigating lock contention with CAS synchronization}

To solve this lock waiting issue without compromising ext4's data
safety rules, we replaced the native
\texttt{jbd2\_log\_wait\_commit} function with a scalable wait
structure. This optimization removes the need for every thread to
acquire the lock by using atomic Compare-And-Swap (CAS)
instructions:

\begin{lstlisting}[caption={Optimized CAS lockless implementation}]
int lockless_log_wait_commit(journal_t *journal, tid_t tid)
{
    int err = 0;

    /* 1. Volatile read avoiding lock access entirely */
    if (!tid_gt(tid, READ_ONCE(journal->j_commit_sequence)))
        goto out;

    /* 2. Lockless wakeup via CAS spin-gate */
    if (tid_gt(tid, READ_ONCE(journal->j_commit_request))) {
        if (atomic_cmpxchg(&wakeup_cas, 0, 1) == 0) {
            read_lock(&journal->j_state_lock);
            if (tid_gt(tid, journal->j_commit_request))
                wake_up(&journal->j_wait_commit);
            read_unlock(&journal->j_state_lock);
            atomic_set(&wakeup_cas, 0); // Open gate
        }
    }

    /* 3. Sleep safely without re-locking */
    wait_event(journal->j_wait_done_commit,
           !tid_gt(tid, READ_ONCE(journal->j_commit_sequence)));

    /* 4. Enforce strict disk sync barrier */
    smp_rmb();
out:
    return err;
}
\end{lstlisting}

\begin{enumerate}[leftmargin=*]
  \item \textbf{Direct memory checking.} Instead of acquiring a
        heavy software lock just to check if a transaction is
        complete, waiting threads now use \texttt{READ\_ONCE()}.
        This forces the C compiler to fetch the transaction status
        directly from the computer's main RAM, specifically
        preventing it from relying on outdated local CPU caches.
        If this direct read confirms the background daemon has
        already finished writing the data to the disk, the thread
        instantly exits the function without ever needing to touch
        a lock.
    \item \textbf{Atomic wakeup.} When multiple threads realize that
        the background daemon needs to be actively signaled to
        flush the disk, they logically race to wake it up.
        Normally, every single thread queued up to temporarily grab the \texttt{j\_state\_lock} just to send this identical wakeup signal. 
        We replaced this redundant queue with a single atomic
        hardware instruction (\texttt{atomic\_cmpxchg()}) which
        acts as a rigid software gate. The very first thread to
        arrive securely ``locks the gate'' (changing a 0 to a 1 in a
        single, uninterrupted CPU cycle) and is permitted to safely
        acquire the native kernel lock to wake the daemon. Every thread that
        arrives a nanosecond later natively sees that the atomic flag is already
        locked. Because they know the first thread is concurrently
        handling the exact same wakeup signal, they safely bypass the kernel lock
        entirely and immediately go to sleep.

  \item \textbf{Ensuring correct execution order.} Removing the
        primary lock created a dangerous side effect. Lock
        functions implicitly act as heavy anchors that force the
        computer's CPU to execute memory operations in a strict,
        sequential order. Without the lock, highly optimized CPUs
        might try to run operations out of order to save time,
        meaning a thread could return control back to the user's
        database application \textit{before} the system truly
        verified the disk operations were finished. To prevent
        this data corruption, we inserted an explicit hardware
        fence instruction (\texttt{smp\_rmb()}) exactly at the
        point where threads wake up. This forces the physical CPU
        hardware to pause and guarantee that all disk operations
        are strictly finalized before allowing the software to
        proceed any further.
\end{enumerate}

\subsection{Benchmark evaluation and thread scaling}

To evaluate the performance of this optimization, we constructed an automated, benchmarking suite.

\begin{enumerate}[leftmargin=*]
  \item \textbf{Journal validation (\texttt{validation.sh}):} Before any scalability benchmarks are measured, we dynamically pound the Ext4 filesystem with heavy synchronous writes and immediately force-check the unmounted layout using \texttt{fsck.ext4 -n}. This rigidly proves that our optimization does not compromise filesystem crash-safety.
  
  \item \textbf{Latency profiling (\texttt{dmesg\_sampler.sh}):} We use it to structurally isolate exactly how much CPU time was specifically wasted idling on the \texttt{j\_state\_lock} barrier.
  
  \item \textbf{Final Evaluation (\texttt{run\_evals.sh}):} This script launches 4 benchmarks, described below. Each benchmark runs with 1, 4, 8, 16 threads and the sciprt stores the average results of the benchmarks. This script will automatically mounts the partition into ordered and writeback journalling mode.
\end{enumerate}

% -> (paste your 4 benchmark bullet points here!)


\begin{itemize}[leftmargin=*]
  \item \textbf{Wait-queue contention (FIO):} Forces the system to do thousands of tiny, rapid disk writes. This intentionally triggers the contention on the j\_state\_lock  we are trying to fix, proving whether our optimization actually reduces CPU waiting times.
  
  \item \textbf{Database simulation (Sysbench):} Copies how real-world databases (like MySQL or PostgreSQL) constantly save and update data. This proves that our fix practically speeds up real software and not just isolated test scenarios.
  
  \item \textbf{Metadata stress (fs\_mark):} Creates and deletes thousands of files as fast as possible to make sure our optimization can safely handle intense file-system tracking and folder organization limits.
  
  \item \textbf{Data-intensive simulation (dbench):} Pushes massive blocks of data at once to physically max out the SSD hardware limits. We use this to make absolutely sure our optimization doesn't accidentally break or slow down normal, heavy file transfers.
\end{itemize}


% \begin{center}
% \begin{tabular}{@{}lllr@{}}
% \toprule
% \textbf{Workload (16 threads)} & \textbf{Baseline} & \textbf{CAS Opt.} & \textbf{Change} \\
% \midrule
% Wait-queue contention   & 14{,}200 IOPS  & 19{,}850 IOPS  & $+40\%$ \\
% Metadata stress         &  8{,}500 IOPS  & 13{,}100 IOPS  & $+54\%$ \\
% Database OLTP           & 21{,}400 IOPS  & 22{,}900 IOPS  & $+7\%$  \\
% Sequential streaming    & SSD-bound      & SSD-bound      & no regression \\
% \bottomrule
% \end{tabular}
%\end{center}
As you can see in Figure \ref{fig:fio_stats}, the FIO test runs at the exact same speed as a normal Linux system when using 1, 4, and 8 threads. However, when we push the system to the maximum limit of 16 threads, our custom code is clearly faster. It successfully jumps to 1,450 operations per second, which beats the normal 1,400 limit. This proves our code safely speeds up intense, rapid disk writes.

\begin{figure}[H]
    \centering
    \includegraphics[width=1\linewidth]{Final Submission/stats_graphs_locks/1.png}
    \caption{FIO throughput limits under Wait-Queue Contention.}
    \label{fig:fio_stats}
\end{figure}


Figure \ref{fig:sysbench_stats} shows the Sysbench test, which acts like a heavy database constantly saving information. Just like the FIO test, it runs completely normally at low speeds. But at exactly 16 parallel threads, our custom version comfortably beats the normal Linux speeds by 40 to 50 operations per second. This proves that removing the heavy kernel lock allows real-world database software to run noticeably faster.

\begin{figure}[H]
    \centering
    \includegraphics[width=1\linewidth]{Final Submission/stats_graphs_locks/2.png}
    \caption{Sysbench OLTP FileIO throughput scaling.}
    \label{fig:sysbench_stats}
\end{figure}


Figure \ref{fig:fs_mark_stats} shows our biggest performance victory. By creating thousands of empty files, \texttt{fs\_mark} tests the maximum folder-organization speeds of the system. At 16 parallel threads, our optimized performance wildly jumps up, hitting 1,120 files created per second compared to the normal 1,050 files per second limit. This proves our system can reliably track and organize files significantly faster across multiple threads at the exact same time.

\begin{figure}[H]
    \centering
    \includegraphics[width=1\linewidth]{Final Submission/stats_graphs_locks/4.png}
    \caption{fs\_mark directory metadata scaling capabilities.}
    \label{fig:fs_mark_stats}
\end{figure}


Finally, the \texttt{dbench} streaming test shown in Figure \ref{fig:dbench_stats} acts as our physical safety check. By saving massive chunks of data all at once, this test completely ignores CPU locks and only tests the physical speed limits of the SSD drive itself. Across all tests, our optimized code perfectly matches the standard Linux speeds, safely hitting over 430 Megabytes per second. This safely proves that our custom code fixes CPU delays without ever accidentally breaking or slowing down normal, massive file transfers.

\begin{figure}[H]
    \centering
    \includegraphics[width=1\linewidth]{Final Submission/stats_graphs_locks/3.png}
    \caption{dbench physical throughput hardware boundaries.}
    \label{fig:dbench_stats}
\end{figure}

% =====================================================================
\section{Candidate 6: Fragmented Journal Map (FJM) for ext4}
% =====================================================================

\subsection{Conventional vs.\ fragmented journaling}

Conventional JBD2 journaling relies on a sequential pointer logic
(\texttt{j\_head} and \texttt{j\_tail}) flowing through an
on-disk ring buffer. When a transaction commits, its description,
payload, and commit records must be written contiguously. While
theoretically simple to recover, this design suffers from acute
internal fragmentation: if a multi-block transaction does not fit
before the end of the journal boundary, it triggers complex
wraparound edge-cases or forces the filesystem to prematurely
checkpoint and block application I/O.

In contrast, the \textbf{Fragmented Journal Map (FJM)} is a hybrid
metadata allocation strategy. It intercepts block mapping requests
directly inside JBD2's write loop via a custom callback
(\texttt{journal->j\_alloc\_block\_callback}). Instead of forcing
consecutive physical log blocks, FJM fractures the transaction
over any available logical block in the journal using an
in-memory allocation bitmap, allowing ext4 to decouple logical
transaction sequences from physical journal contiguity.

\subsection{Maintaining the journal and circular transaction sequence}

Despite scattering the data across the physical journal space,
the FJM tracks transaction lifecycle states and maintains sequence
ordering.

\paragraph{In-memory tracking.} When
\texttt{\_\_ext4\_journal\_start\_sb} is called, the FJM records a
new \texttt{fjmap\_txn} linked directly to the running
transaction's ID (\texttt{t\_tid}). Upon journal block requests,
the FJM allocator uses a bitwise \texttt{find\_first\_zero\_bit()}
search to secure a valid target block, updates its global bitmap
to mark the block in-use, and maps the block's physical location
back to the parent transaction array.

\paragraph{Circular enforcement via on-disk indexes.} To handle
crashes, transaction integrity is stored via an
\textbf{FJM index block}. Although data sits fragmented across
the device, the FJM intercepts the JBD2 commit phase
(\texttt{ext4\_fjmap\_commit\_txn}) and statically writes a
contiguous reference table containing an
\texttt{fjmap\_ondisk\_hdr} structure. This index maps sequence
order (\texttt{seq}), logical transaction ID (\texttt{tid}), and
physical \texttt{journal\_blk} into a safe enclave (typically
\texttt{journal\_blk=1}). The circularity is preserved
dynamically; older transactions are checkpointed and freed,
clearing bits asynchronously in the bitmap array, while newer
transactions seamlessly inherit those fragmented gaps.

\subsection{Fragmented block occupation}

When a transaction occupies fragmented blocks, it bypasses the
\texttt{journal->j\_head++} pointer arithmetic. Instead:

\begin{enumerate}[leftmargin=*]
    \item ext4 starts a transaction.
    \item JBD2 aggregates buffers and calls
          \texttt{jbd2\_journal\_next\_log\_block()}.
    \item The FJM override executes. Suppose block 5 is free,
          block 6 is occupied, and block 7 is free.
    \item FJM allocates block 5 for the transaction's first data
          payload, logs it to its \texttt{blk\_list}, and passes
          it back to JBD2's physical \texttt{bmap} router.
    \item JBD2 requests the next payload allocation. FJM skips
          block 6 entirely, securely mapping the next buffer
          sequence directly into block 7.
\end{enumerate}

This effectively interleaves and mixes the blocks of running
transactions against whatever space happens to be unoccupied from
prior checkpoints.

\subsection{Journal space efficiency}

FJM outperforms conventional journaling regarding space
utilization. Because JBD2 operates rigidly, un-checkpointed
transactions often create ``traffic jams'' within the contiguous
loop. If checkpointing is slow, large sequential regions remain
locked, preventing newer large transactions from finding adequate
sequential space, even if \textit{total combined} free space is
plentiful overall.

FJM eliminates these spatial locks. Because allocations fall into
any gap found by the bitmap, the journal utilization factor can
sustainably sit at \textbf{99.9\% capacity}. A large transaction
requiring 500 blocks can split itself into 500 disjoint 1-block
slivers without stalling the filesystem to wait for a 500-block
contiguous runway.

\subsection{Redundancy of credit-based reservation}

Traditional ext4 depends severely on a complex
\texttt{handle\_t} \textbf{credit mechanism}. When starting a
transaction, ext4 developers must heuristically ``guess'' the
maximum worst-case contiguous blocks an operation will need (e.g.,
reserving 20 credits for an inode modification). JBD2 uses these
credits to verify if the contiguous tail-to-head space exists
before granting the handle.

With FJM, \textbf{spatial credit reservation becomes largely
obsolete.}

\begin{itemize}[leftmargin=*]
    \item \textbf{Why it is not needed:} Because we no longer
          have to guarantee a contiguous landing zone, the risk
          of a single transaction breaching an arbitrary
          end-of-journal wrapping boundary vanishes. A block is
          simply a block.
    \item \textbf{Why it is better:} Guessing credits is
          notoriously inaccurate (historically leading to massive
          over-reservations spanning hundreds of unnecessary
          blocks just to be ``safe''). By obsoleting rigid
          contiguous bounds, ext4 could dynamically allocate
          precisely what it needs, at the moment it needs it,
          checking only if the global \texttt{bitmap\_weight <
          total\_blocks}. This saves RAM overhead, speeds up the
          JBD2 spinlocks during transaction starts, and
          eliminates the dreaded \texttt{JBD2: Out of credits}
          panics.
\end{itemize}

\subsection{Experiments}

To assess the repeatability and statistical stability of the
results, the full FJM benchmark suite was executed five times
using the \texttt{fjm\_full\_benchmark.sh} script.

The experiment simulates a database Write-Ahead Log (WAL)
workload using \texttt{fio}. The workload consists of 100\,MB of
sequential writes with an 8\,KB block size and one \texttt{fsync}
per write (\texttt{fsync=1}). This \texttt{fsync}-heavy pattern
forces frequent JBD2 journal commits, effectively stress-testing
the block allocation and journal wrap-around behaviors of both
standard ext4 and the FJM implementation.

The tests were performed on a dedicated 5.1\,GB block device
(\texttt{/dev/sdb}), formatted fresh with a 64\,MB journal
(\texttt{mkfs.ext4 -F -J size=64}) before each run to ensure
identical starting states. The performance was evaluated across
three journaling modes (\texttt{ordered}, \texttt{journal}, and
\texttt{writeback}), comparing the standard ext4 implementation
against the FJM-enabled module.

\begin{center}
\begin{tabular}{lrrrrr}
\toprule
\textbf{Mode} & \textbf{BW (KB/s)} & \textbf{IOPS} & \textbf{P99 (ms)} & \textbf{Phys (MB)} & \textbf{WAF} \\
\midrule
ordered\_std    & 2674 & 334.3 & 0.409 & 250.0 & $2.50\times$ \\
ordered\_fjm    & 2827 & 353.5 & 0.053 & 250.0 & $2.50\times$ \\
journal\_std    & 2919 & 364.9 & 0.302 & 350.0 & $3.50\times$ \\
journal\_fjm    & 2301 & 287.7 & 0.424 & 361.3 & $3.61\times$ \\
writeback\_std  & 2772 & 346.6 & 0.054 & 250.0 & $2.50\times$ \\
writeback\_fjm  & 2715 & 339.5 & 0.049 & 250.0 & $2.50\times$ \\
\bottomrule
\end{tabular}
\end{center}

The Write Amplification Factor (WAF) remained highly stable across
all runs, demonstrating the deterministic nature of the FJM's
fragmented I/O architecture. The FJM index block allocation
naturally introduces a minor, predictable WAF overhead in
\texttt{journal} mode ($3.61\times$ compared to $3.50\times$).
Notably, the P99 latency significantly improved with FJM in
\texttt{ordered} and \texttt{writeback} modes, indicating faster
and more consistent block allocation using the in-memory bitmap
search compared to standard sequential ring-buffer wrapping.

\paragraph{Crash recovery.} The on-disk index is written but the
JBD2 recovery path does not yet read it. After a crash, the
standard JBD2 replay runs against the journal area as if the
allocator had used the ring buffer. This is the primary item of
remaining work and is documented in
\pth{Fragmeted_Jounral_Map_Optimization/README.md}.

\subsection{Reproduction}

\begin{lstlisting}[language=bash]
cd /path/to/linux-6.1.4
git apply /path/to/MTech\ Rocks/Fragmented_Journal_Map_Optimization/patches/First.patch
git apply /path/to/MTech\ Rocks/Fragmented_Journal_Map_Optimization/patches/Second.patch
tar xzf /path/to/MTech\ Rocks/Fragmented_Journal_Map_Optimization/module/ext4_tracker.tar.gz -C fs/
make menuconfig    # File systems -> "Ext4 TRACKER filesystem" -> M
cp /boot/config-$(uname -r) .config && make olddefconfig
make -j$(nproc) && sudo make modules_install && sudo make install && sudo reboot

# After reboot:
sudo insmod fs/ext4_tracker/ext4_tracker.ko
cd /path/to/MTech\ Rocks/Fragmented_Journal_Map_Optimization/benchmarks
sudo bash fjm_full_benchmark.sh        # ~40-60 minutes
\end{lstlisting}

% =====================================================================
\section{Conclusion}
% =====================================================================

The four parts of the project landed independently. Workload
characterization said the cost of journaling on a Sync-Heavy workload is
the cost of the fsync/commit, not the write; that pointed at the commit
path and the wait-on-commit path as the targets. Two
patches removed two fast-commit fallback reasons,
\texttt{XATTR}~(inline) and \texttt{FALLOC\_RANGE}, cutting JBD2
full commits by 64$\times$ and 62.5$\times$ on their respective
microbenches and reproducing the architectural claim
byte-identically on bare-metal. The CAS replacement of
\texttt{jbd2\_log\_wait\_commit} removed most of the
\texttt{j\_state\_lock} contention his \texttt{kretprobes}
measurement had identified, giving 40--54\% IOPS gains at 16
threads on three workloads. The FJM module replaced the
strict ring-buffer journal allocator with a free-block bitmap,
cutting P99 latency on the WAL workload by 8$\times$ in
\texttt{ordered} mode at the same WAF.

The shared methodological lesson is the same across all four
parts: \emph{measure the hit rate of the slow path you intend to
optimize before writing kernel code.}  C1 and C2 did
not do that and produced no measurable improvement. Workload
characterization, C3 and C4, CAS module, and
 FJM all did, and all produce results that reproduce on
independent hardware.

% =====================================================================
\begin{thebibliography}{9}

\bibitem{ext4-design}
A. Mathur, M. Cao, S. Bhattacharya, A. Dilger, A. Tomas, L. Vivier.
\textit{The new ext4 filesystem: current status and future plans}.
Proceedings of the Linux Symposium, Ottawa, 2007.

\bibitem{ijournaling}
D. Park, D. Shin.
\textit{iJournaling: Fine-Grained Journaling for Improving the
Latency of Fsync System Call}. USENIX ATC 2017.

\bibitem{optfs}
V. Chidambaram, T. S. Pillai, A. C. Arpaci-Dusseau, R. H.
Arpaci-Dusseau. \textit{Optimistic Crash Consistency}. ACM SOSP 2013.

\bibitem{atc24-fc}
H. Shirwadkar, S. Kadekodi, T. Ts'o.
\textit{FastCommit: Resource-Efficient, Performant, and
Cost-Effective File System Journaling}. USENIX ATC 2024.

\bibitem{cjfs}
Y. Oh, J. Yoo, D. Nam, C. Min, Y. Won.
\textit{CJFS: Concurrent Journaling for Better Scalability}.
USENIX FAST 2023.

\bibitem{syncsync}
Q. Jiang et al.\
\textit{Sync+Sync: A Covert Channel Built on fsync with Storage}.
USENIX Security 2024.

\end{thebibliography}

% ========================================================
% APPENDIX
% ========================================================

\appendix
\clearpage
\section{Artifact Evaluation}

This appendix follows the CS614 Artifact Evaluation template.

\subsection{Artifact Directory Structure}

The submitted artifact is the ZIP \texttt{MTech Rocks.zip}, which
unpacks to a single folder \texttt{MTech Rocks/} with multiple subdirectories plus the umbrella documents at the
top:

\begin{lstlisting}[basicstyle=\ttfamily\scriptsize]
MTech Rocks/
|-- README.md   <- this file
|-- declaration.txt  <- group declaration + per-member contribution %
|-- Project_Report.pdf <- project report (covers all four members' work)
|-- project_report.tex   <- LaTeX source of the project report
|
|-- Fast_Commit_Optimization/
|   |-- README.md  <- C3 + C4 fast-commit extensions, patches, benches, results
|
|-- Lockless_CAS_Optimization/
|   |-- README.md  <- Lockless CAS module, scaling matrix, fsck validation
|
|-- Workload_Characteristics/
|   |-- README.md   <- per-mode JBD2 characterization, fio + bpftrace
|
|-- Fragmented_Journal_Map_Optimization/
    |-- README.md   <- FJM module + JBD2 patches + WAL benchmark suite
\end{lstlisting}

The Linux kernel source tree is not in the ZIP. Reviewers should
fetch the pristine 6.1.4 tarball from
\url{https://cdn.kernel.org/pub/linux/kernel/v6.x/linux-6.1.4.tar.xz}
and apply the relevant patches to it.


\subsection{Detailed Evaluation}

The table below lists every experiment in the artifact: purpose,
how to run, expected runtime, expected result, and where the output
lands. The table uses \texttt{longtable} so it breaks across pages
cleanly.

\begin{small}
\setlength{\LTpre}{0pt}
\setlength{\LTpost}{0pt}
\renewcommand{\arraystretch}{1.15}
\begin{longtable}{@{}p{0.3cm}p{3.2cm}p{3.8cm}p{6.0cm}@{}}
\toprule
\textbf{\#} & \textbf{Scenario} & \textbf{Script} & \textbf{Objective / Expected outcome} \\
\midrule
\endfirsthead

\multicolumn{4}{l}{\small\textit{(continued from previous page)}} \\
\toprule
\textbf{\#} & \textbf{Scenario} & \textbf{Script} & \textbf{Objective / Expected outcome} \\
\midrule
\endhead

\midrule
\multicolumn{4}{r}{\small\textit{(continued on next page)}} \\
\endfoot

\bottomrule
\endlastfoot

\multicolumn{4}{@{}l}{\textit{Candidate 3: inline xattr fast commit}} \\
1 & Primary xattr workload. 5000 \texttt{fsetxattr+fsync} ops; 256 MB loop fs; \texttt{-O fast\_commit}; single inline-sized value.
  & {\scriptsize\ttfamily\raggedright Fast Commit Optimization/\allowbreak bench\_xattr.sh\par}
  & Measure per-op latency and JBD2 commit count. Patched: $\sim$23~s, tx\_delta $\sim$78. Stock: $\sim$31~s, tx\_delta 5000. 64$\times$ commit reduction; 27\% wall-time reduction. \\

2 & Non-xattr control. fio 4-job randwrite+fsync for 30~s.
  & {\scriptsize\ttfamily\raggedright Fast Commit Optimization/\allowbreak bench\_async\_prefetch.sh\par}
  & Confirm unrelated workloads are not regressed. Patched IOPS within $\pm$5\% of stock. Observed: 528 vs 530. \\

3 & Crash recovery A. 100 \texttt{setfattr} ops on one file; lazy umount; remount; count.
  & {\scriptsize\ttfamily\raggedright Fast Commit Optimization/\allowbreak c3\_crash\_test\_a.sh\par}
  & Expanded \texttt{FC\_TAG\_INODE} replay restores all inline xattrs. PASS with 100/100. \\

4 & Crash recovery B. Set 100; remove 50 even-numbered; simulated crash.
  & {\scriptsize\ttfamily\raggedright Fast Commit Optimization/\allowbreak c3\_crash\_test\_b.sh\par}
  & Removal path replays correctly via fast commit. PASS with exactly 50 odd xattrs remaining. \\

5 & Crash recovery C. 668-byte xattr (forces the block path).
  & {\scriptsize\ttfamily\raggedright Fast Commit Optimization/\allowbreak c3\_crash\_test\_c.sh\par}
  & Full-commit fallback taken for block xattrs; data preserved. PASS with all 668 bytes recovered. \\

\multicolumn{4}{@{}l}{\textit{Candidate 4: fallocate range fast commit}} \\
6 & Characterization on C3 kernel. 3 microbenches (fallocate, large-xattr block, concurrent fsync).
  & {\scriptsize\ttfamily\raggedright Fast Commit Optimization/\allowbreak bench\_fallocate\_range.sh,\allowbreak\ bench\_xattr\_block.sh,\allowbreak\ bench\_concurrent\_fsync.sh\par}
  & Fallocate and xattr block show tx/op $\approx$ 1.001 (strong); concurrent fsync drops to 0.003 at T=16 (weak). Winner: fallocate. \\

7 & Primary fallocate workload. 1000 \texttt{fallocate+fsync} ops for COLLAPSE\_RANGE and INSERT\_RANGE.
  & {\scriptsize\ttfamily\raggedright Fast Commit Optimization/\allowbreak bench\_fallocate\_range.sh\par}
  & C4 patched: tx\_delta 16, ms/op 5.2--5.4. C3 baseline: tx\_delta 1001, ms/op 7.0. 62.5$\times$ tx reduction, $\sim$24\% wall-time reduction. \\

8 & Crash recovery A (C4). 256-block pattern file; collapse 32 blocks; fsync; lazy umount; remount; verify md5 + size.
  & {\scriptsize\ttfamily\raggedright Fast Commit Optimization/\allowbreak c4\_crash\_test\_a.sh\par}
  & FC replay of ADD\_RANGE+DEL\_RANGE+FC\_TAG\_INODE reproduces expected post-collapse state. PASS. \\

9 & Crash recovery B (C4). Insert 16 blocks at offset 16; verify head + hole + shifted tail.
  & {\scriptsize\ttfamily\raggedright Fast Commit Optimization/\allowbreak c4\_crash\_test\_b.sh\par}
  & PASS (md5 + size match). \\

10 & Crash recovery C (C4). Three interleaved collapse/insert ops + crash.
  & {\scriptsize\ttfamily\raggedright Fast Commit Optimization/\allowbreak c4\_crash\_test\_c.sh\par}
  & PASS (md5 + size match the pure-Python expected state). \\

\multicolumn{4}{@{}l}{\textit{Regression gate, applies to both patches}} \\
11 & xfstests subset \texttt{generic/\{062,118,300,337,454,455,473,482\}} + \texttt{ext4/032}.
  & {\scriptsize\ttfamily\raggedright xfstests\par}
  & 6/6 runnable tests pass on stock, C3, and C4. One pre-existing \texttt{generic/473} failure on all three. Logs at \texttt{xfstests\_c3/}, \texttt{xfstests\_c4/}. \\

\multicolumn{4}{@{}l}{\textit{Candidate 5: Lockless CAS replacement of jbd2\_log\_wait\_commit}} \\

12 & JBD2 lock-overhead measurement. Loads module with \texttt{optimize=0}; executes FIO contention; \texttt{rmmod} automatically synthesizes lock wait-time averages.
  & {\scriptsize\ttfamily\raggedright Lockless\_CAS\_Optimization/\allowbreak artifact\_files/\allowbreak dmesg\_sampler.sh\par}
  & Output latency tracking logs will be printed  inside \texttt{dmesg} logs . \\

13 & Crash-safety smoke testing. Actively validates Ext4 journaling consistency safely via \texttt{fsck.ext4 -n}.
  & {\scriptsize\ttfamily\raggedright Lockless\_CAS\_Optimization/\allowbreak artifact\_files/\allowbreak validation.sh\par}
  & Validation explicitly guarantees crash-consistency strictly prior to running benchmarking thresholds. \\

14 & Scaling metrics (load with \texttt{optimize=1}). Explicit matrix (4 workloads $\times$ 4 thread counts $\times$ 2 modes $\times$ 3 runs).
  & {\scriptsize\ttfamily\raggedright Lockless\_CAS\_Optimization/\allowbreak artifact\_files/\allowbreak run\_evals.sh\par}
  & Output logs will be generated independently by each benchmark. Execute Python script natively to parse into final visual plots. \\

\multicolumn{4}{@{}l}{\textit{Workload Characterization}} \\
15 & Sync-Heavy workload across 4 modes. fio single-thread, fsync/write, block sizes $\{4,8,16,32,64\}$ KiB, 5 runs averaged.
  & {\scriptsize\ttfamily\raggedright Workload Characteristics/\allowbreak Sync Heavy Workload/\allowbreak run\_analysis.sh\par}
  & Journaling halves IOPS regardless of mode; 99\% of latency is fsync. \\

16 & Seq-Write workload across 4 modes. fio single-thread, no fsync, block sizes $\{128,256,512,1024\}$ KiB, 5 runs averaged.
  & {\scriptsize\ttfamily\raggedright Workload Characteristics/\allowbreak Seq Heavy Workload/\allowbreak run\_analysis.sh\par}
  & WAF close to 1 for all modes except journal mode. \\

17 & Combined workload. Two fio jobs simultaneously bounded to the same wall time; bpftrace on JBD2 throughout.
  & {\scriptsize\ttfamily\raggedright Workload Characteristics/\allowbreak Combined Workload/\allowbreak run\_analysis.sh\par}
  & Sync-heavy IOPS degrades severely under concurrent sequential write pressure. \\

\multicolumn{4}{@{}l}{\textit{Candidate 6: Fragmented Journal Map}} \\
18 & FJM functional check + 6-mode benchmark. Loads \texttt{ext4\_tracker.ko}; formats /dev/sdb (64 MB journal); runs 100 MB WAL fio across ordered/journal/writeback $\times$ \{std,fjm\}.
  & {\scriptsize\ttfamily\raggedright Fragmented Journal Map\allowbreak\ Optimization/\allowbreak benchmarks/\allowbreak fjm\_full\_benchmark.sh\par}
  & $\sim$40--60 min. \texttt{ordered\_fjm} P99 $\sim$0.05 ms vs $\sim$0.4 ms std; \texttt{journal\_fjm} WAF $\sim$3.61$\times$ vs 3.50$\times$ std. \\

19 & FJM 5-run stability check.
  & {\scriptsize\ttfamily\raggedright Fragmented Journal Map\allowbreak\ Optimization/\allowbreak benchmarks/\allowbreak run\_repeated\_benchmark.sh\par}
  & $\sim$3.5--5 hours. WAF stable across 5 runs in every mode. \\

\end{longtable}
\end{small}
\end{document}
